\documentclass[11pt]{article}

\usepackage[margin=1in]{geometry}
\usepackage{amsmath,amssymb,amsfonts,amsthm}
\usepackage{array}
\usepackage{booktabs}
\usepackage{graphicx}
\usepackage{hyperref}
\usepackage{url}
\usepackage{xcolor}

\hypersetup{
  colorlinks=true,
  linkcolor=blue,
  citecolor=blue,
  urlcolor=blue
}

\newtheorem{proposition}{Proposition}
\newtheorem{theorem}{Theorem}
\newtheorem{lemma}{Lemma}
\newtheorem{assumption}{Assumption}
\newtheorem{corollary}{Corollary}

\newcommand{\method}{Gauss6/FullVA}
\newcommand{\tfe}{TFE}
\newcommand{\artifact}[1]{\texttt{#1}}
\newcommand{\figpath}{figures}
\newcolumntype{P}[1]{>{\raggedright\arraybackslash}p{#1}}
\newcommand{\plotfigure}[3]{%
  \begin{figure}[t]
    \centering
    \IfFileExists{#1}{%
      \includegraphics[width=0.95\linewidth]{#1}%
    }{%
      \fbox{\parbox{0.9\linewidth}{Missing figure: \nolinkurl{#1}}}%
    }
    \caption{#2}
    \label{#3}
  \end{figure}
}

\title{Lie-Group FullVA Integration for Cylindrical Lower-Pair Mechanisms:\\
A Conditional Sixth-Order Formal-Order Comparison with Separated Four-Example Evidence}
\author{Jingquan Wang}
\date{v047 artifact-backed paper build}

\begin{document}
\maketitle

\noindent\textbf{Status note (2026-09-17).} This file is a legacy internal
report kept for the record. The current manuscript is \texttt{main\_cmame.tex}
in the same directory. It replaces the proof route sketched here by the exact
stage identity: the implemented stage system is three-stage Gauss collocation
of the reduced joint-coordinate equations written in absolute coordinates, so
the stage residual at the lifted Gauss stage is zero and the order theorem
rests on the P1, P2, and P6 interfaces only. The observed orders above six
quoted below are finite-window slopes of a sixth-order collocation method and
are not order claims.

\begin{abstract}
This paper reports the v047 evidence for a Lie-group FullVA integrator applied
to cylindrical lower-pair mechanisms.  The accepted \method{} path is evaluated
on the cylindrical-chain benchmark and on four ASME-style mechanism examples
with separated evidence roles: single pendulum, double pendulum, four link,
and slider crank.  On the smooth
cylindrical h-sweep, the method records observed position and velocity orders
7.161 and 7.066.  The sharp Brown--McPhee friction case is order-reduced on the
coarse h-sweep, but high-order behavior reappears under deeper fixed
refinement, leaving a practical coarse-regime caveat rather than a closed-form
correctness failure.

The main claim of this paper is therefore a conditional formal-order
comparison: within the current artifact scope, the accepted \method{} path
has a conditional sixth-order Lie-group FullVA theorem on the selected smooth
branch, while the four examples provide separated method-side evidence roles.
The single- and double-pendulum rows support dynamic-order evidence; the
four-link and slider-crank rows support lower-pair mechanism coverage.  The
local paper-style $m=3$ Lobatto-\tfe{} formula target is order five.  A
full-\tfe{} replacement would strengthen the reproduction claim, but it is not
needed for this bounded formal-order comparison.

The paper also separates this accepted result from a stricter reproduction
claim.  The original Chaturvedi--Sandu--Sandu paper has its own full \tfe{}
formulation.  In this v047 implementation, ``full-\tfe{} replacement'' means
replacing the accepted Gauss6/FullVA stage residual by paper-derived
temporal-finite-element weak rows inside the nonlinear solve.  That replacement
gate remains open and machine-readable as
\nolinkurl{full_tfe_stage_replacement=false}.  The open replacement gate is
reported as a limitation and as future reproduction work; it does not weaken
the accepted conditional method/proof claim or the separated four-example
evidence roles.
\end{abstract}

\section*{Reviewer-Facing Summary}

The current paper should be read as an artifact-backed status paper, not as a
claim that every research gate is closed.  The accepted method claim is the
conditional sixth-order \method{} path under the direct residual-bridge proof
route.  The cylindrical-chain benchmark records smooth position/velocity
orders 7.161/7.066.  The four ASME-style examples have separated roles:
single and double pendula provide dynamic-order evidence, while four link and
slider crank provide mechanism-coverage and reaction/constraint diagnostics.

The paper-\tfe{} claim is separate.  The encoded paper formula map uses
$m=3$ Gauss--Lobatto temporal finite elements and has expected order five in
the formula-mapping audit, but this is not yet an accepted method replacement.
The machine-readable gate remains
\nolinkurl{full_tfe_stage_replacement=false}.  The remaining missing object is
an eight-row source-free lower-pair closure that replaces the terminal bridge
inside the 132-row nonlinear stage residual while preserving the paper-\tfe{}
order target.
Equivalently, the source paper is not missing a \tfe{} method; the missing
piece is our artifact-backed substitution of that paper-style residual into
the v047 solver.
The newest endpoint-pose trajectory smoke has narrowed that missing object:
the terminal-equivalent endpoint row closes velocity but loses order, while a
nearby nonterminal row preserves neither terminal closure nor order.  The next
repair therefore has to be a nonterminal endpoint-pose predictor that closes
terminal velocity on trajectory without collapsing the h-sweep order.  Sliding
the extrapolation points almost onto the endpoint closes terminal velocity, but
it reproduces the same 4.142/2.305 smooth-order loss; quadratic nonterminal
extrapolation does the same.

Routine paper validation should use the read-only validators, not the full
generator.  The historical \artifact{run\_v047.py} file is an accumulated
audit harness containing positive evidence, negative screens, and artifact
generation code; the latest recorded full generation time is
$2386.76\,\mathrm{s}$.  It is not needed to check the accepted four-example
claim or rebuild this PDF.

\section{Scope and Claims}

The v047 pipeline has one primary paper goal and one optional stronger
reproduction goal.  The primary goal is to validate a conditional sixth-order
Lie-group multibody integrator for cylindrical lower-pair mechanisms.  This is the
accepted \method{} path: it has a conditional sixth-order theorem on the
selected smooth branch by the direct 132-row residual bridge, records smooth
position/velocity orders 7.161/7.066 in the current h-sweep, and has
four-example coverage with separated evidence roles.  The single- and
double-pendulum rows provide dynamic-order evidence, while the four-link and
slider-crank rows provide closed-loop constraint/reaction checks that are not
interpreted as accepted asymptotic order estimates.  The
optional stronger goal audits whether a
complete paper-derived temporal finite-element stage residual can replace the
current stage equations without relying on endpoint projection, endpoint
source data, or terminal-row substitution.  This second goal is intentionally
stricter and is not needed for the formal-order comparison claim.

The source paper itself is therefore not being accused of lacking a full
\tfe{} formulation.  The open v047 item is narrower and implementation-facing:
we have not yet produced an accepted nonlinear solve in which the paper-style
\tfe{} stage rows replace the Gauss6/FullVA stage rows while retaining the
required order, rank, residual, and endpoint-closure evidence.

The full-\tfe{} replacement is therefore not needed to state the four-example
method evidence or the bounded formal-order comparison.  It is needed only for the
stronger statement that the implemented nonlinear residual is the original
paper-style temporal finite-element method rather than the accepted
conditional Gauss6/FullVA method proved through the direct 132-row residual
bridge with separate diagnostic closure evidence.
Keeping this boundary explicit is important because otherwise the formal-order
comparison could be confused with a solved paper-method reproduction.

The distinction matters for interpreting the numerical order.  The accepted
\method{} path is high order on smooth data.  The current source-free
final-stage \tfe{} closure candidate and the newer endpoint-pose velocity
predictor smoke are not accepted replacements: they close or nearly close
terminal velocity, but they do not preserve the target trajectory order.  The
paper therefore reports the \method{} theorem and method evidence as the
positive result and the full \tfe{} replacement as an open proof and
implementation gate.

\subsection{Terminology and Claim Boundary}

The phrase \emph{independent full-\tfe{} stage replacement} is used in a narrow
technical sense.  It does not mean merely adding an endpoint correction to an
otherwise accepted Gauss trajectory.  It means replacing the stage residual
rows themselves by paper-derived temporal finite-element rows and solving those
rows inside the nonlinear trajectory problem.

\begin{table}[t]
\centering
\caption{Terms used in the v047 paper and pipeline.}
\label{tab:terms}
\begin{tabular}{p{0.27\linewidth}p{0.65\linewidth}}
\toprule
Term & Meaning in this paper \\
\midrule
\method{} &
The accepted Gauss collocation/FullVA trajectory residual: a conditional
order-six method path on the selected smooth branch by the direct 132-row
residual bridge. \\
Paper-\tfe{} &
The temporal finite-element formula family encoded from the paper coefficient
maps; the v047 formula-mapping audit uses $m=3$ Gauss--Lobatto nodes with
expected order five. \\
Independent full-\tfe{} replacement &
All 132 stage rows are replaced by paper-derived stage rows and are solved
inside Newton, with accepted h-sweep order, residual, rank, and endpoint checks. \\
Source-free lower-pair closure &
The lower-pair closure rows do not receive an external endpoint source and do
not rely on endpoint-boundary source data. \\
Terminal-row replacement &
A diagnostic shortcut that inserts raw endpoint velocity rows in place of
paper rows.  This is useful for debugging but is not accepted as full \tfe{}. \\
Projection &
A post-step velocity correction.  Projection is accepted only on the historical
Gauss diagnostic path, not on the independent full-\tfe{} replacement path. \\
\bottomrule
\end{tabular}
\end{table}

Thus the current open issue is not whether the four examples run.  They do.
It is whether the stricter paper-\tfe{} residual can replace the Gauss stage
equations without losing the order and endpoint behavior already demonstrated
by the accepted \method{} path.  The occasional spelling ``FTE'' in notes should
be read as \tfe{}, temporal finite element.

\subsection{Claim Map}

The current result has one primary accepted formal-order comparison, one accepted
four-example method gate, and one optional open reproduction claim.  The
primary claim is that \method{} is a conditional sixth-order method in the
artifact scope used here: it has observed smooth orders 7.161/7.066, whereas
the encoded local paper-style $m=3$ Lobatto-\tfe{} target has expected order
five.  The open claim is narrower: a paper-derived
full-\tfe{} stage residual has not yet independently replaced the Gauss stage
equations.

\begin{table}[t]
\centering
\caption{Current claim boundary.  The order labels are intentionally separated
because they refer to different residuals.}
\label{tab:claim-map}
\begin{tabular}{P{0.29\linewidth}P{0.17\linewidth}P{0.43\linewidth}}
\toprule
Claim & Status & Evidence or blocker \\
\midrule
Formal-order comparison relative to the local paper-style \tfe{} target &
Accepted bounded comparison &
\method{} has a conditional order-six theorem by the direct residual-bridge
proof and records smooth observed orders 7.161/7.066; the encoded $m=3$
Gauss--Lobatto \tfe{} formula target has expected order five.  This is not a
full residual reproduction claim. \\
Gauss6/FullVA cylindrical-chain method &
Accepted method path &
Direct 132-row bridge: 96 non-dynamic rows supply row-local \(O(h^7)\)
residual input and 36 Newton--Euler rows vanish by direct substitution on the
lifted Gauss stage, giving the conditional local \(O(h^7)\)/global \(O(h^6)\)
claim under the stated interfaces. \\
Four ASME-style method rows &
Accepted method gate &
Single and double pendulum method-side FullVA rows are accepted; four-link and
slider-crank closed-loop kinematic FullVA plus reaction-dynamics rows are
accepted. \\
Paper formula mapping &
Mapped diagnostic &
The encoded $m=3$ Gauss--Lobatto \tfe{} operator has expected order five, but
mapping the formulas is not the same as solving a full replacement residual. \\
Source-free terminal lower-pair closure &
Useful but not accepted as full \tfe{} &
Raw terminal endpoint velocity closes to $4.01\times10^{-16}$, but smooth
position order is 3.523. \\
Independent full-\tfe{} stage replacement &
Open &
The missing eight source-free lower-pair closure rows must preserve terminal
closure and restore the paper-\tfe{} h-sweep order without endpoint source,
projection, terminal-row replacement, or a target-direction oracle. \\
\bottomrule
\end{tabular}
\end{table}

\subsection{Comparator Definition}

The comparison in Proposition~1 is deliberately local and reproducible.  The
comparator is the paper-style $m=3$ Gauss--Lobatto \tfe{} formula target
encoded in the v047 formula-mapping audit, not an independently accepted
implementation of every source-paper residual row.  This keeps the accepted
claim aligned with the available evidence: the pipeline can prove that the
conditional \method{} path has formal order six while the local paper-style
target has expected order five, while still reporting that the complete source-paper
residual reproduction remains open.

\begin{table}[t]
\centering
\caption{Comparator definition for the paper-facing formal-order comparison.}
\label{tab:comparator-definition}
\begin{tabular}{P{0.24\linewidth}P{0.24\linewidth}P{0.40\linewidth}}
\toprule
Object & Role in this paper & Current evidence boundary \\
\midrule
Accepted \method{} path &
Method being claimed &
Gauss/FullVA nonlinear residual with a conditional order-six theorem by direct
132-row residual bridge; smooth observed orders 7.161/7.066; four-example
evidence separated into dynamic-order pendulum rows and residual/reaction
coverage on four-link/slider crank. \\
Local paper-style $m=3$ Lobatto-\tfe{} target &
Comparator used for Proposition~\ref{prop:formal-order-comparison} &
Formula-mapping target with expected order five.  It is sufficient for the
current formal-order comparison because the comparison is an order and
evidence-scope claim. \\
Complete source-paper \tfe{} residual &
Not used as the accepted comparator &
The source paper has a full \tfe{} formulation, but v047 has not yet accepted
an implementation in which all paper-derived stage rows replace the
Gauss6/FullVA rows inside Newton.  This is the open
\nolinkurl{full_tfe_stage_replacement=false} reproduction gate. \\
v046/rA baseline &
Historical reproducibility anchor &
Useful for repository regression context, but not the main paper-facing
comparator for Proposition~\ref{prop:formal-order-comparison}. \\
\bottomrule
\end{tabular}
\end{table}

\begin{table}[t]
\centering
\caption{Original-paper boundary versus accepted v047 claim.}
\label{tab:source-paper-v047-boundary}
\begin{tabular}{P{0.24\linewidth}P{0.32\linewidth}P{0.32\linewidth}}
\toprule
Question & Source-paper/local target & Accepted v047 claim \\
\midrule
What is being compared? &
The local paper-style $m=3$ Lobatto-\tfe{} formula target, whose expected
order is five in the formula-mapping audit. &
The accepted \method{} path, a conditional order-six Gauss/FullVA one-step map
with smooth observed orders 7.161/7.066. \\
What is implemented as the method? &
The complete source-paper \tfe{} residual is not used as the accepted
comparator because the independent replacement gate remains open. &
The nonlinear solve uses the accepted 132-row Gauss/FullVA residual and the
audited projection/KKT endpoint closure. \\
What is proved by the artifacts? &
They prove the local order-five paper-style target boundary, not a completed
source-paper residual reproduction. &
They support a paper-facing formal-order comparison inside the current
v047 artifact scope and separated four-example ASME evidence gate. \\
What remains outside the claim? &
Full source-paper residual reproduction requires accepted paper-derived stage
rows inside Newton. &
\nolinkurl{full_tfe_stage_replacement=false} remains an explicit caveat and
future stronger gate. \\
\bottomrule
\end{tabular}
\end{table}

\begin{table}[t]
\centering
\caption{Acceptance criteria for the paper-facing formal-order comparison.}
\label{tab:formal-order-acceptance}
\begin{tabular}{P{0.25\linewidth}P{0.48\linewidth}P{0.15\linewidth}}
\toprule
Criterion & Required evidence & Status \\
\midrule
Accepted method path &
The claimed method is the \method{} one-step map in
Section~\ref{sec:accepted-one-step-integrator}, not a diagnostic
paper-\tfe{} substitution. &
Pass \\
Order exceedance &
The direct residual-bridge theorem gives order six under the stated interfaces,
and the accepted smooth h-sweep records position/velocity orders 7.161/7.066,
above the encoded local $m=3$ Lobatto-\tfe{} expected order-five target. &
Pass \\
Four-example gate &
The accepted method rows cover single pendulum, double pendulum, four link,
and slider crank under the same read-only ASME gate. &
Pass \\
Comparator boundary &
The comparison is against the local paper-style formula target; complete
source-paper residual reproduction remains a separate gate. &
Pass \\
Caveat separation &
Sparse speed, sharp-friction coarse-regime behavior, and
\nolinkurl{full_tfe_stage_replacement=false} are explicit caveats and do not
contradict the narrower comparative claim. &
Pass \\
Read-only reproducibility &
The paper-claim, paper-package, four-ASME, full-\tfe{} gap, repair-spec, and
pipeline-output validators check the claim without invoking the full v047
generator. &
Pass \\
\bottomrule
\end{tabular}
\end{table}

\subsection{Main Comparative Proposition}

\begin{proposition}[Formal-order comparison]
\label{prop:formal-order-comparison}
Within the current v047 artifact scope, the accepted \method{} path supports a
conditional formal-order comparison relative to the encoded local paper-style
$m=3$ Lobatto-\tfe{} target.
\end{proposition}

\begin{proof}[Evidence and proof sketch]
The accepted method path is the Gauss/FullVA residual on the cylindrical
lower-pair benchmark.  Its order statement is the conditional theorem proved
by the direct 132-row residual bridge: on the selected smooth compact branch,
the 96 non-dynamic rows supply row-local \(O(h^7)\) residual input, the 36
Newton--Euler rows vanish by direct substitution on the lifted Gauss stage,
and endpoint/Newton perturbations are retained at \(O(h^7)\) before the
local-to-global \(O(h^6)\) transfer.  The smooth h-sweep records observed
position and velocity orders 7.161/7.066 as downstream numerical evidence, and
the same pipeline accepts the four-example evidence gate.  Only the single-
and double-pendulum dynamic rows are used as method-order evidence; the
four-link and slider-crank rows remain closed-loop constraint/reaction checks
whose near-roundoff trajectory differences are not interpreted as convergence
order.  The comparator used for this proposition is not a complete
source-paper residual implementation; it is the locally encoded $m=3$
Gauss--Lobatto \tfe{} formula target, whose expected order is five in the
formula-mapping audit.  Therefore the artifact-backed comparison is between a
conditional order-six Gauss/FullVA method path and a local order-five
paper-style target.  This proves the bounded paper-facing comparative
statement claimed here.

The proposition is intentionally not a full residual reproduction theorem.  A
full source-paper reproduction statement would require the still-open
independent full-\tfe{} stage replacement gate, where the paper-derived rows
replace the Gauss6/FullVA stage residual inside the nonlinear solve and pass
the rank, residual, endpoint-closure, and h-sweep checks without endpoint
source data, terminal-row replacement, projection, or a target-direction
oracle.  That stronger statement remains outside the accepted
formal-order comparison.
\end{proof}

\subsection{Contributions}

The paper is organized around artifact-backed claims rather than a
single final theorem.  The current contributions are:
\begin{itemize}
  \item a conditional order-six integrator comparison for the cylindrical
  lower-pair benchmark: the accepted \method{} path is proved through the
  direct 132-row residual bridge and records smooth position and velocity
  orders 7.161/7.066, while the local paper-style $m=3$ Lobatto-\tfe{} formula
  target has expected order five;
  \item accepted method-side evidence for four ASME-style examples, including
  lower-pair graph, rank, closed-loop kinematic FullVA, and reaction-dynamics
  checks for the closed-loop examples;
  \item a paper-\tfe{} formula-mapping and substitution audit that keeps the
  expected fifth-order Lobatto-\tfe{} claim distinct from the accepted Gauss6
  method claim;
  \item a source-free final-stage lower-pair closure that removes endpoint
  source data, terminal-row replacement, and projection while closing raw
  terminal velocity to roundoff;
  \item a negative but useful compression-audit sequence showing which
  source-free eight-row closure families do not yet provide an independent
  full-\tfe{} replacement.
\end{itemize}

\subsection{Role of the Full-TFE Replacement}

The full-\tfe{} replacement is not needed to support the accepted Gauss6
result: the order-six claim is conditional on the retained compact-branch,
endpoint, binding, and solver-scale interfaces, with smooth numerical evidence
downstream.  The replacement is optional future work for a stronger
source-paper reproduction claim.  If the paper is to claim a temporal
finite-element stage method rather than a Gauss-stage method with endpoint
diagnostics, then the rows solved by Newton must be the paper-derived weak/TFE
rows themselves.  In particular, the lower-pair endpoint closure must be
produced by stage-local source-free rows, not by an endpoint source, a
final-row selector, a projection step, or a local target-Jacobian oracle.

This distinction is why the pipeline keeps running targeted compression audits
after the four-example method gate is accepted.  Those audits are not redoing
the four examples.  They are isolating the specific eight-row lower-pair
closure law that would make the paper-\tfe{} residual independently viable.

\subsection{Reader Status Checklist}

The current status can be read from four separate questions.  This separation
prevents the accepted method evidence from being confused with the still-open
paper-\tfe{} replacement evidence.

\begin{table}[t]
\centering
\caption{Status checklist for interpreting the v047 claims.}
\label{tab:reader-status}
\begin{tabular}{P{0.31\linewidth}P{0.17\linewidth}P{0.40\linewidth}}
\toprule
Question & Answer & Evidence boundary \\
\midrule
Is the main formal-order comparison accepted? &
Yes &
\method{} is the accepted conditional order-six method path; the smooth
h-sweep records 7.161/7.066, above the local paper-style $m=3$
Lobatto-\tfe{} expected order-five target. \\
Does the four-example method-evidence gate pass? &
Yes &
The accepted method gate covers single pendulum and double pendulum as
dynamic-order evidence, and four link and slider crank as lower-pair
constraint/reaction coverage on the \method{} path. \\
Is the production method sixth order? &
Conditionally &
The theorem proves local \(O(h^7)\)/global \(O(h^6)\) on the selected smooth
branch under the compact-tube, endpoint, implementation-binding, and
solver-scale interfaces; the smooth h-sweep records observed orders
7.161/7.066. \\
Is the paper $m=3$ Lobatto-\tfe{} formula mapped? &
Yes &
The formula-mapping audit records expected order five, but this is not by
itself a solved method replacement. \\
Is the independent full-\tfe{} stage replacement accepted? &
No &
The source-free terminal closure closes endpoint velocity, but the replacement
candidate remains order-limited and the machine-readable gate remains
\nolinkurl{full_tfe_stage_replacement=false}. \\
\bottomrule
\end{tabular}
\end{table}

\section{Numerical Model}

The v047 cylindrical-chain benchmark replaces twist-locked prismatic joints by
cylindrical joints.  The cylindrical lower pair constrains axis alignment and a
perpendicular point relation while leaving axial translation and relative spin
free.  The main residual has dimension 132.  Its sparse Jacobian pattern has
2637 nonzeros; column coloring uses 90 colors, and the row-VJP backend uses 60
row colors while preserving trajectory agreement to roundoff.

The accepted trajectory integrator uses a three-stage Gauss collocation
structure with FullVA variables and is claimed through the conditional
order-six residual-bridge theorem.  Each body carries position, orientation,
velocity, angular velocity, acceleration, and angular acceleration information
inside the stage solve.  Lower-pair constraints enter at position, velocity,
and acceleration levels through the stage residual.  Endpoint velocity closure
is audited through both the historical projection path and a square endpoint
KKT residual.

At a schematic level, the accepted path solves
\begin{equation}
  R_{\mathrm{Gauss6}}(q_n,v_n,a_n,Z;h)=0,
\end{equation}
where $Z$ denotes the FullVA stage variables.  Lower-pair constraints appear as
\begin{equation}
  \Phi(q)=0,\qquad
  \Phi_q(q)v=0,\qquad
  \Phi_q(q)a+\dot{\Phi}_q(q,v)v=0.
\end{equation}
The paper-\tfe{} replacement target instead seeks a stage residual
\begin{equation}
  R_{\tfe{}}(q_n,v_n,a_n,Z;h)=
  \begin{bmatrix}
  R_x & R_y & R_z & R_{\mathrm{NE}} & R_{\Phi,v} & R_{\Phi,a}
  \end{bmatrix}^{T}=0,
\end{equation}
with the same 132-row stage budget, but with rows derived from the paper
temporal finite-element construction rather than from the Gauss stage formula.
This replacement is counted as full only when the resulting trajectory solve,
not just a one-step tangent or residual evaluation, passes the acceptance
checks.

\subsection{Accepted One-Step Integrator}
\label{sec:accepted-one-step-integrator}

The integrator claimed in Proposition~\ref{prop:formal-order-comparison} is the
one-step map
\begin{equation}
  x_{n+1}=\Psi_h^{\mathrm{G6FVA}}(x_n),
  \qquad x_n=(q_n,v_n,a_n),
\end{equation}
defined by the accepted Gauss6/FullVA residual and its recorded endpoint
closure audit.  It is not the complete paper-\tfe{} replacement residual.
For each step, v047 uses the following evidence-backed construction:
\begin{enumerate}
  \item assemble the cylindrical lower-pair constraints and FullVA stage
  unknowns $Z$ for the current state $x_n$ and step size $h$;
  \item solve the 132-row nonlinear stage system
  $R_{\mathrm{Gauss6}}(q_n,v_n,a_n,Z;h)=0$ to the Newton tolerance recorded in
  the CSV and JSON artifacts;
  \item recover the raw endpoint by the Gauss collocation quadrature and the
  local Lie-group retraction used by the v047 trajectory code;
  \item audit endpoint velocity closure through the square endpoint KKT path
  while keeping historical projection residual/correction measurements as
  diagnostics;
  \item emit the per-step residual, constraint, order, runtime, and plotting
  artifacts consumed by the paper and pipeline validators.
\end{enumerate}
The theorem endpoint object is the velocity-level/KKT closure functional.
It is evaluated on the accepted Gauss6/FullVA path.  Raw projection residuals
are diagnostics, not theorem inputs or full-\tfe{} replacements.  This is the
integrator whose smooth cylindrical-chain h-sweep records 7.161/7.066
position/velocity orders and whose four-example method gate is accepted.

\subsection{Order Numbers}

There are two order targets in the current evidence.  The production method is
the conditional order-six Gauss6/FullVA path, and the smooth benchmark measures
position/velocity orders 7.161/7.066 on the tested range.  The paper formula
mapping audit separately encodes an $m=3$ Gauss--Lobatto \tfe{} operator with
expected order five.  An earlier source-free terminal-closed candidate has
smooth orders 3.523/4.828.  The latest endpoint-pose velocity predictor
trajectory smoke closes terminal velocity even more tightly on the terminal
row, but gives only 4.142/2.305 position/velocity order on the three-$h$
smooth refinement.  These are useful closure diagnostics, not accepted
replacements, because they still degrade the trajectory order.

\begin{table}[t]
\centering
\caption{Order conventions used in the paper.  The rows are deliberately
separated because they refer to different nonlinear residuals.}
\label{tab:order-conventions}
\begin{tabular}{P{0.28\linewidth}P{0.24\linewidth}P{0.36\linewidth}}
\toprule
Path & Order statement & Interpretation \\
\midrule
Accepted \method{} path &
Conditional order-six method; observed smooth orders 7.161/7.066 &
This is the method used for the accepted cylindrical-chain and four-example
evidence claims. \\
Paper $m=3$ Lobatto-\tfe{} formula map &
Expected order five &
This is a formula-mapping diagnostic.  It becomes a method claim only after the
paper-derived rows replace the stage residual inside Newton. \\
Current source-free terminal-closed \tfe{} candidate &
Smooth orders 3.523/4.828 &
Terminal velocity is closed without endpoint source, terminal-row replacement,
or projection, but order is too low for an accepted full-\tfe{} replacement. \\
Endpoint-pose velocity predictor smoke &
Smooth orders 4.142/2.305 &
The terminal-equivalent \nolinkurl{0p00} row closes terminal velocity to
$8.124\times10^{-17}$, but the trajectory order remains too low; the nearby
nonterminal \nolinkurl{0p01} row leaves terminal velocity at
$6.255\times10^{-6}$. \\
\bottomrule
\end{tabular}
\end{table}

\section{Proof Status}

The proof ledger separates the load-bearing theorem from generated numerical
evidence.  The accepted \method{} path is not certified merely by a generic
``Gauss collocation inherits order'' slogan.  The load-bearing argument fixes
the implemented 132-row FullVA residual, inserts the classical three-stage
Gauss local defect into that same stage object, uses the 96-row non-dynamic
certificate and the 36-row Newton--Euler direct-substitution identity, and then
passes the resulting \(O(h^7)\) stage residual through the compact-tube
Kantorovich, endpoint, inexact-Newton, and local-to-global transfer lemmas.
Sparse AD and row/column coloring do not change this order because they only
change how the same Newton residual and Jacobian are evaluated.

\begin{assumption}[Smooth accepted-path order assumptions]
\label{ass:g6fva-order}
For the smooth cylindrical benchmark and reported h-sweep, assume:
\begin{enumerate}
  \item the square FullVA residual is locally equivalent to a regular smooth
  constrained flow in a local Lie-group chart;
  \item the three-stage Gauss collocation equations are solved to a Newton
  tolerance below the sixth-order truncation scale;
  \item the local chart and retraction used for endpoint reconstruction are
  smooth diffeomorphisms on the tested trajectory tube;
  \item the theorem endpoint output is the velocity-level/KKT closure
  functional on the same compact branch, while raw endpoint projection
  monitors remain diagnostics and do not supply the endpoint theorem input;
  \item the sparse backend evaluates the same residual and Jacobian as the
  dense backend.
\end{enumerate}
\end{assumption}

\begin{theorem}[Order of the accepted \method{} map]
\label{thm:g6fullva-order}
Under Assumption~\ref{ass:g6fva-order}, interpreted as the retained
accepted-branch P1--P3/P6 compact-tube interfaces with the P4/P5 direct
\(96+36\) residual bridge as the only consumed residual-value certificate, the
accepted \method{} one-step map in
Section~\ref{sec:accepted-one-step-integrator} has local defect \(O(h^7)\) and
reduced-chart/reported position--velocity grid error \(O(h^6)\) on fixed smooth
time intervals, for reported grids whose transitions satisfy the predeclared
branch, endpoint, proof-norm, and \(\eta_h^{\rm tube}\le c_\eta h^7\)
solver-envelope conditions.  For driven or otherwise nonautonomous rows, the
residual \(F_{A,h,t_n}\), stage times \(t_n+c_i h\), endpoint map, and proof
norm are fixed on the same transition slice before constants or diagnostics
are used.  This is a conditional order-six statement for
the accepted branch-selected map only.  It retains P6 as a theorem-level
solver-envelope condition, does not prove a P7 residual-to-error transfer, does not
prove residual, reaction, multiplier, or constraint-output error estimates,
does not use source-policy rows, and does not use or complete the separate
primitive/Taylor T3 route.  Equivalently, the theorem asserts only the
local defect $O(h^7)$ and same-initial-state global error $O(h^6)$ under these
retained interfaces and only for this displayed reduced/reporting grid scope.
\end{theorem}

\begin{proof}
By Assumption~\ref{ass:g6fva-order}, the proof is carried out on a fixed
smooth compact branch in the local Lie-group chart.  The classical
three-stage Gauss local defect is inserted into the implemented FullVA stage
object, not into a different reduced model.  On that branch, the 96
non-dynamic rows provide the row-local \(O(h^7)\) residual input and the 36
Newton--Euler rows vanish by direct substitution on the lifted Gauss stage.
Thus the accepted 132-row bridge supplies the \(O(h^7)\) stage residual
defect used by the conditional theorem.  Endpoint closure and the inexact
Newton iterate add only \(O(h^7)\) perturbations under the retained
compact-tube hypotheses, and the local-to-global transfer then gives the
\(O(h^6)\) grid estimate for the same branch-selected reported grid.  The
proof order is the constrained Lie-group BDF order of obligations, not an
estimate transfer: fix the discrete object, insert the smooth constrained
branch, prove the local defect for that fixed object, and only then apply the
stability recurrence and reporting map.  No BDF local-error constant,
hidden-constraint estimate, multiplier-history estimate, or multistep
zero-stability theorem is imported into this one-step FullVA proof.
The solver residual bound is a retained P6 envelope fixed before the run, not a
conclusion drawn from fixed-tolerance logs, and the residual-to-error P7
interface remains open.  Sparse and dense backends evaluate the same
residual/Jacobian path, so backend choice is an implementation diagnostic, not
an additional order argument.  Residual-only mechanism rows, source-policy
rows, and the separate primitive/Taylor T3 route are not theorem inputs and cannot
be combined with the direct residual bridge to strengthen the theorem.
\end{proof}

\begin{lemma}[Order of the local paper-style \tfe{} target]
\label{lem:tfe-target-order}
The proof-level order target for the $m=3$ Gauss--Lobatto paper-style
\tfe{} formula used in this manuscript is five.
\end{lemma}

\begin{proof}
The source-paper family is compared through the temporal finite-element
formula-order convention $p_{\mathrm{TFE}}(m)=2m-1$.  Substituting $m=3$
gives $p_{\mathrm{TFE}}(3)=5$.  This is distinct from the source paper's
reported DAE pendulum order reduction.
\end{proof}

\begin{theorem}[Accepted formal-order comparison theorem]
\label{thm:accepted-formal-order-comparison}
Assume Assumption~\ref{ass:g6fva-order}.  Also assume the current v047 artifact contract:
the read-only validators accept the four-example coverage rows, the single-
and double-pendulum dynamic rows provide the method-side order evidence, the
four-link and slider-crank rows are treated only as closed-loop
constraint/reaction checks, the smooth cylindrical h-sweep records the reported
7.161/7.066 position/velocity orders, the local paper-style $m=3$
Lobatto-\tfe{} formula target has expected order five, and the claim boundary records
\nolinkurl{full_tfe_stage_replacement=false}.  Then the v047 evidence supports
the following and only the following paper-facing statement: the accepted
\method{} one-step map is a conditional sixth-order method relative to the
local order-five paper-style \tfe{} target and is covered by the four
ASME-style examples under the stated evidence boundary.  The theorem does not assert complete source-paper
\tfe{} residual reproduction.
\end{theorem}

\begin{proof}
\emph{Order.}  Theorem~\ref{thm:g6fullva-order} gives the conditional
order-six result for the accepted \method{} one-step map through the direct
132-row residual bridge under the smooth compact-branch and solver-scale
interfaces.  The generated convergence artifacts give the independent
numerical check used by the paper: the smooth cylindrical h-sweep records
position and velocity orders 7.161/7.066.

\emph{Four-example evidence.}  The same artifact contract accepts the
single pendulum, double pendulum, four link, and slider crank method rows on
the \method{} path.  These rows are accepted as method evidence rows, not as
full-\tfe{} replacement rows.

\emph{Comparison.}  The comparator in this paper is the locally encoded
paper-style $m=3$ Lobatto-\tfe{} formula target.  Lemma~\ref{lem:tfe-target-order}
fixes its expected order as five. Combining this order-five comparator with
the conditional order-six accepted \method{} map proves the formal-order
comparison statement inside the stated
artifact scope.

\emph{Boundary.}  A complete source-paper residual reproduction would require
accepted paper-derived stage rows inside the 132-row nonlinear solve, without
endpoint source data, terminal-row replacement, projection, or a
target-direction oracle.  The current artifact contract records
\nolinkurl{full_tfe_stage_replacement=false}; therefore that stronger
statement is excluded rather than silently assumed.  The theorem proves only
the accepted formal-order comparison claim.
\end{proof}

\begin{corollary}[Artifact-backed formal-order comparison acceptance]
\label{cor:artifact-backed-acceptance}
Under Assumption~\ref{ass:g6fva-order} and the current read-only v047 artifact
checks, Proposition~\ref{prop:formal-order-comparison} is accepted within the stated
artifact scope.
\end{corollary}

\begin{proof}
Theorem~\ref{thm:g6fullva-order} gives the conditional order-six direct
residual-bridge theorem for the accepted \method{} one-step map, and
Theorem~\ref{thm:accepted-formal-order-comparison} packages the order,
four-example, comparator, and boundary evidence into the formal paper-facing
claim.  The generated convergence artifacts record smooth position/velocity
orders 7.161/7.066, and the minimal four-example validator accepts single
pendulum, double pendulum, four link, and slider crank on the accepted method
path.  The comparator definition keeps the local $m=3$ Lobatto-\tfe{}
expected-order-five target separate from the complete source-paper residual
reproduction gate.  The full-\tfe{} gap and repair-spec validators keep
\nolinkurl{full_tfe_stage_replacement=false} explicit, so the corollary
accepts only the formal-order comparison and not the stronger reproduction
claim.
\end{proof}

This is deliberately not a proof of the independent paper-\tfe{} replacement.
The current paper-\tfe{} rows are either formula maps, local tangent probes, or
diagnostic substitutions unless they replace the stage residual inside the
nonlinear h-sweep and pass the acceptance test in
Table~\ref{tab:full-tfe-acceptance-test}.  The velocity-level/KKT closure
functional is the theorem endpoint object on the accepted Gauss path; raw
endpoint projection measurements are diagnostics and do not discharge the
endpoint raw-defect hypothesis or define an accepted full-\tfe{} method.

\begin{table}[t]
\centering
\caption{Proof/status boundary inherited from the order-proof ledger.}
\label{tab:proof-status}
\begin{tabular}{P{0.29\linewidth}P{0.22\linewidth}P{0.37\linewidth}}
\toprule
Object & Status & Reason \\
\midrule
\method{} smooth cylindrical method &
Conditional order-six argument &
Direct 132-row bridge gives local \(O(h^7)\) residual defect and global
\(O(h^6)\) grid error under the retained compact-branch, endpoint, and
solver-scale interfaces. \\
Sparse Jacobian backend &
Order-preserving backend change &
Exact AD and sparse linear algebra evaluate the same residual/Jacobian to the
same Newton tolerance. \\
Endpoint closure path &
Diagnostic caveat &
Raw endpoint/projection residuals remain diagnostics; the theorem endpoint
object is the velocity-level/KKT closure functional, so projection monitors are
not used to claim independent full-\tfe{} acceptance. \\
Paper $m=3$ Lobatto-\tfe{} mapping &
Mapped but not accepted as a method &
Expected order five is recorded for the formula map, but the rows have not yet
produced an accepted full replacement h-sweep. \\
Independent full-\tfe{} replacement &
Open proof obligation &
The source-free terminal closure is algebraically useful, but the current
h-sweep remains order-limited and the gate is still
\nolinkurl{full_tfe_stage_replacement=false}. \\
\bottomrule
\end{tabular}
\end{table}

\section{Validation Protocol}

Every reported method claim is tied to generated artifacts in the v047 results
directory.  The core h-sweeps use
three step sizes $h\in\{0.04,0.02,0.01\}$ against an $h=0.005$ reference unless
otherwise noted.  The four ASME-style examples are:
\begin{itemize}
  \item single pendulum;
  \item double pendulum;
  \item four link;
  \item slider crank.
\end{itemize}

The validation records CSV rows, summary JSON entries, plots, and a generated
report.  Accepted method rows require small residuals, clean constraints, and
consistent observed order or a documented reference policy.  Diagnostic rows
are kept separate from accepted rows when they use projection, endpoint-source
data, local tangent probes, terminal-row replacement, or incomplete \tfe{}
stage substitution.

\subsection{Proposition Evidence Traceability}

Table~\ref{tab:proposition-traceability} records the evidence chain for
Proposition~\ref{prop:formal-order-comparison}.  It is deliberately read-only: every
entry is checked from existing CSV, JSON, PDF, and ledger artifacts rather than
by rerunning the full v047 generator.

\begin{table}[t]
\centering
\caption{Traceability for the formal-order comparison proposition.}
\label{tab:proposition-traceability}
\begin{tabular}{P{0.26\linewidth}P{0.43\linewidth}P{0.19\linewidth}}
\toprule
Proposition component & Artifact evidence & Read-only check \\
\midrule
Accepted \method{} one-step map &
Summary JSON, convergence CSV, and
Section~\ref{sec:accepted-one-step-integrator} &
Paper-claim validator \\
Smooth high-order behavior &
Convergence CSV and convergence figure record 7.161/7.066 orders &
v047 and paper-package validators \\
Four-example acceptance &
Summary JSON ASME gate with single pendulum, double pendulum,
four link, and slider crank &
Minimal four-ASME validator \\
Comparator boundary &
Comparator table and paper claim ledger separate the local
$m=3$ Lobatto-\tfe{} target from complete source-paper residual reproduction &
Paper-claim validator \\
Open full-\tfe{} caveat &
Full-\tfe{} gap ledger and repair specification keep
\nolinkurl{full_tfe_stage_replacement=false} explicit &
Full-\tfe{} gap and repair-spec validators \\
Corollary acceptance &
Assumption--lemma--theorem--corollary chain accepts the proposition only inside
the current read-only v047 artifact scope &
Paper-claim validator \\
Repository synchronization &
Order-proof ledger, pipeline audit, and version ledger carry the same claim
boundary &
Pipeline-output validator \\
\bottomrule
\end{tabular}
\end{table}

\section{Accepted Gauss6/FullVA Results}

\plotfigure{\figpath/convergence.png}
{Cylindrical-chain convergence for the accepted \method{} path.  The smooth
case is high order; the sharp-friction case is order-reduced on the coarse
h-sweep.}
{fig:convergence}

\begin{table}[t]
\centering
\caption{Cylindrical-chain observed orders for the accepted \method{} path.}
\label{tab:main-order}
\begin{tabular}{P{0.39\linewidth}P{0.18\linewidth}P{0.18\linewidth}P{0.10\linewidth}}
\toprule
Case & Step sizes & Position order & Velocity order \\
\midrule
Smooth, projected velocity & 0.04, 0.02, 0.01 & 7.161 & 7.066 \\
Smooth, raw endpoint       & 0.04, 0.02, 0.01 & 7.161 & 7.066 \\
Sharp, projected velocity  & 0.04, 0.02, 0.01 & 1.894 & 2.683 \\
Sharp, raw endpoint        & 0.04, 0.02, 0.01 & 1.894 & 2.683 \\
\bottomrule
\end{tabular}
\end{table}

Table~\ref{tab:main-order} gives the headline convergence behavior.  The
smooth branch is above sixth order on the tested range.  The sharp branch is
not treated as a silent failure; it is tracked by a Brown--McPhee smoothness
sweep and by deeper fixed refinements.

\plotfigure{\figpath/friction_smoothness_sweep.png}
{Brown--McPhee smoothness sweep.  Reducing the stribeck velocity moves the
problem from a smooth high-order branch to a coarse sharp-friction
order-reduced branch.}
{fig:smoothness}

\begin{table}[t]
\centering
\caption{Sharp-friction refinement evidence.}
\label{tab:sharp}
\begin{tabular}{lccc}
\toprule
Audit & Step window & Position order & Velocity order \\
\midrule
Fixed refinement & 0.02, 0.01, 0.005 & 3.294 & 5.867 \\
Deep fixed refinement & 0.01, 0.005, 0.0025 & 6.398 & 5.870 \\
Ultra fixed refinement & 0.005, 0.0025, 0.00125 & 7.521 & 5.725 \\
\bottomrule
\end{tabular}
\end{table}

The ultra-refinement audit recovers high observed order, but only at increased
work.  Relative to the $h=0.01$ sharp baseline, the high-order window around
$h=0.00125$ costs about $7.25\times$ runtime and $8\times$ steps/Newton
iterations while reducing velocity error by approximately $1.06\times10^6$.
The remaining sharp-friction caveat is therefore a practical coarse-regime
cost issue rather than a complete asymptotic failure.

\section{Four-Example Mechanism Gate}

\plotfigure{\figpath/asme_lower_pair_graph_bridge.png}
{Lower-pair graph bridge over the four ASME-style examples.}
{fig:graph-bridge}

\plotfigure{\figpath/asme_closed_loop_kinematic_fullva.png}
{Closed-loop kinematic FullVA solves for the four-link and slider-crank
examples.}
{fig:closed-loop}

\begin{table}[t]
\centering
\caption{Four-example evidence summary.}
\label{tab:asme}
\begin{tabular}{lp{0.66\linewidth}}
\toprule
Example & Accepted evidence \\
\midrule
Single pendulum &
Exact driven kinematic mapping; full absolute-coordinate driven FullVA
residual; embedded scalar FullVA bridge; reaction-multiplier reconstruction. \\
Double pendulum &
Method-side FullVA mapping with accepted nested local v029/FullVA
self-reference; v046 retained as reproducibility and reference-floor
diagnostic. \\
Four link &
Lower-pair graph bridge, full-row-rank $\Phi_q$ audit, accepted local
closed-loop kinematic FullVA solve, and accepted reaction-dynamics rows. \\
Slider crank &
Lower-pair graph bridge, full-row-rank $\Phi_q$ audit, accepted local
closed-loop kinematic FullVA solve, and accepted reaction-dynamics rows. \\
\bottomrule
\end{tabular}
\end{table}

\begin{table}[t]
\centering
\caption{Quantitative evidence behind the accepted four-example gate.}
\label{tab:asme-quant}
\begin{tabular}{P{0.20\linewidth}P{0.39\linewidth}P{0.29\linewidth}}
\toprule
Example & Quantitative evidence & Artifact path \\
\midrule
Single pendulum &
Absolute-coordinate driven FullVA residual has minimum observed order 6.024
across position, velocity, orientation, and angular velocity; max stage
residual is $3.664\times10^{-12}$. &
Single absolute FullVA runs CSV \\
Double pendulum &
Accepted local v029/FullVA self-reference has minimum orientation/angular
velocity order 6.089; endpoint position/velocity constraint norms stay below
$7.034\times10^{-13}$ and $8.323\times10^{-12}$. &
Double method runs CSV \\
Four link &
Closed-loop kinematic FullVA solve over $h=\{0.02,0.01,0.005\}$ keeps
position/velocity/acceleration constraint norms below
$1.052\times10^{-14}$; reaction-dynamics residual is at most
$1.338\times10^{-13}$. &
Closed-loop kinematic and reaction CSVs \\
Slider crank &
Closed-loop kinematic FullVA solve over $h=\{0.02,0.01,0.005\}$ keeps
position/velocity/acceleration constraint norms below
$1.204\times10^{-15}$; reaction-dynamics residual is at most
$6.492\times10^{-15}$. &
Closed-loop kinematic and reaction CSVs \\
\bottomrule
\end{tabular}
\end{table}

The four-example gate is accepted as coverage, not as four independent
dynamic-order proofs.  This does not mean all
research caveats are closed: sparse speed, the independent full-\tfe{} stage
replacement, and practical sharp-friction coarse-step behavior remain open
items.  The machine-readable summary status for this claim is
\nolinkurl{four_asme_method_rows_accepted_projection_sharp_sparse_caveats}.

\section{Endpoint Closure and TFE Replacement Audit}

\plotfigure{\figpath/endpoint_kkt_closure.png}
{Endpoint KKT residual closure.  The endpoint velocity constraint is enforced
inside a square residual in this diagnostic path.}
{fig:kkt}

The endpoint KKT diagnostic closes endpoint velocity to roundoff while
preserving the accepted \method{} stage equations.  It is important evidence,
but it is not the same as a complete paper-\tfe{} stage replacement.  The full
\tfe{} replacement target is stricter: all 132 stage rows must be replaced by
independently derived paper-\tfe{} rows, and the resulting nonlinear h-sweep
must pass residual, rank, endpoint, and observed-order checks without relying
on endpoint projection, endpoint source data, terminal-row substitution, or
local tangent-only arguments.

\begin{table}[t]
\centering
\small
\setlength{\tabcolsep}{3pt}
\caption{Representative paper-\tfe{} replacement audits.  ``Full'' denotes
accepted independent full-stage replacement, not merely a successful local
solve.}
\label{tab:tfe-audits}
\begin{tabular}{P{0.41\linewidth}P{0.18\linewidth}P{0.18\linewidth}P{0.08\linewidth}}
\toprule
Audit & Max residual & Smooth orders & Full \tfe{}? \\
\midrule
Position-row substitution & $5.68\times10^{-12}$ & 1.227 / 0.731 & No \\
Kinematic-row substitution & $9.57\times10^{-12}$ & 1.406 / 1.542 & No \\
All-row substitution & $9.64\times10^{-12}$ & 1.406 / 1.542 & No \\
Consistent-z0 terminal output & $9.79\times10^{-12}$ & 4.046 / 4.420 & No \\
Lower-pair acceleration terminal output & $8.35\times10^{-12}$ & 4.937 / 4.799 & No \\
Source-free mean-velocity closure & $8.13\times10^{-12}$ & 3.512 / 4.448 & No \\
Source-free final-stage velocity closure & $9.49\times10^{-12}$ & 3.523 / 4.828 & No \\
Source-free order/closure blend, best-order beta & $8.90\times10^{-12}$ & 5.341 / 6.825 & No \\
Endpoint-pose predictor, \nolinkurl{0p00} three-$h$ smoke & $9.76\times10^{-12}$ & 4.142 / 2.305 & No \\
\bottomrule
\end{tabular}
\end{table}

\subsection{Latest Source-Free Final-Stage Closure}

\plotfigure{\figpath/source_free_final_stage_velocity_closure.png}
{Source-free final-stage lower-pair velocity closure.  This audit closes raw
terminal endpoint velocity but remains smooth-position-order limited.}
{fig:final-stage}

The latest lower-pair audit inserts 16 centered acceleration
source-consistency rows and 8 final paper-stage velocity rows.  It removes the
external endpoint source, removes endpoint-boundary source data from the
trajectory, avoids terminal-row replacement, and does not use velocity
projection.  The nonlinear h-sweep solves six smooth/sharp trajectory rows with
maximum residual $9.49\times10^{-12}$ and raw terminal endpoint velocity
$4.01\times10^{-16}$.

This is a real improvement: the terminal-velocity gate is closed by source-free
stage rows.  It is still not an accepted full-\tfe{} replacement because the
smooth position order is 3.523, below the acceptance floor used by the
pipeline.  The current blocker has shifted from terminal closure to preserving
paper-\tfe{} order after the closure rows are inserted into the nonlinear
trajectory solve.

\subsection{Latest Endpoint-Pose Trajectory Smoke}

The latest bounded JSON-only follow-up promotes the local endpoint-pose row
into the nonlinear 132-row residual without running the full artifact
generator.  The target
\nolinkurl{lower_pair_endpoint_pose_velocity_predictor_h_sweep_smoke} keeps
rank 132 and converges all requested smooth short-trajectory steps.  With the
terminal-equivalent \nolinkurl{stage02_convex_pose_velocity_0p00_z} row, the
three-$h$ refinement over $h=[0.04,0.02,0.01]$ against reference
$h=0.005$ closes raw terminal endpoint velocity to
$8.124\times10^{-17}$, but the position/velocity orders are only 4.142/2.305.

This resolves an ambiguity in the local row-space evidence.  The endpoint row
is not merely untested at trajectory level: it closes terminal velocity inside
Newton, but loses the paper-\tfe{} order.  Moving slightly away from the
endpoint does not fix the problem.  The nearby nonterminal
\nolinkurl{stage02_convex_pose_velocity_0p01_z} row converges, but leaves
terminal endpoint velocity at $6.255\times10^{-6}$ and gives only 2.345/1.878
orders.  The remaining repair target is therefore a nonterminal endpoint-pose
predictor that keeps terminal velocity closed on trajectory while recovering
the paper-\tfe{} order target.

The best non-projection follow-up in this family is the terminal-limit
extrapolated row
\nolinkurl{stage02_convex_pose_velocity_extrapolate_0p01_0p02_z}.  The default
two-$h$ smoke has \nolinkurl{projection_used=false}, reduces terminal velocity
to $1.254\times10^{-7}$, and records 6.412/3.775 orders.  However, the
three-$h$ smooth refinement gives only 4.054/2.364 position/velocity order and
still has maximum terminal velocity $1.254\times10^{-7}$, with the finest-$h$
row at $1.244\times10^{-8}$.  It is therefore a useful trajectory near-miss,
not an accepted replacement.

Moving the extrapolation points closer to the terminal endpoint improves only
the terminal velocity.  The default two-$h$ sweep over
\nolinkurl{0p005/0p01}, \nolinkurl{0p002/0p005}, and \nolinkurl{0p001/0p002}
has \nolinkurl{projection_used=false}, terminal velocities
$3.135\times10^{-8}$, $6.270\times10^{-9}$, and $1.254\times10^{-9}$, and
orders 6.621/4.648, 6.599/4.566, and 6.591/4.521.  A three-$h$ refinement of
\nolinkurl{stage02_convex_pose_velocity_extrapolate_0p001_0p002_z} gives
terminal velocities $1.254\times10^{-9}$ overall and $1.244\times10^{-10}$ at
the finest step, but only 4.143/2.307 order.  The ultra-near
\nolinkurl{stage02_convex_pose_velocity_extrapolate_0p00001_0p00002_z} row
finally satisfies the terminal-velocity tolerance without projection, with max
terminal velocity $1.254\times10^{-13}$ and finest-$h$ terminal velocity
$1.254\times10^{-14}$, but its three-$h$ order is still 4.142/2.305 and
\nolinkurl{smooth_order_ok=false}.  This is terminal-equivalent limiting
evidence, not a full-\tfe{} replacement.

A three-point quadratic extrapolation check reaches the same conclusion without
evaluating the terminal row.  The
\nolinkurl{stage02_convex_pose_velocity_quadextrapolate_*} laws estimate the
weight-zero closure from three strictly nonterminal rows and keep
\nolinkurl{projection_used=false}.  The default two-$h$ sweep over
\nolinkurl{0p01/0p02/0p05}, \nolinkurl{0p005/0p01/0p02}, and
\nolinkurl{0p002/0p005/0p01} records terminal velocities
$2.029\times10^{-11}$, $2.029\times10^{-12}$, and $2.029\times10^{-13}$, but
all three velocity orders remain about 4.508.  The best row,
\nolinkurl{stage02_convex_pose_velocity_quadextrapolate_0p002_0p005_0p01_z},
closes the terminal-velocity tolerance on the three-$h$ refinement, with max
terminal velocity $2.029\times10^{-13}$ and finest-$h$ terminal velocity
$8.254\times10^{-15}$, but its order is still 4.142/2.305.  Thus quadratic
nonterminal extrapolation is also terminal-equivalent limiting evidence rather
than a full-\tfe{} replacement.

A projection-like terminal-tangent follow-up rules out the simplest correction
of that nonterminal row.  The full
\nolinkurl{stage02_convex_pose_velocity_0p01_terminalproj_z} correction fails
in the reference run with residual divergence $6.739\times10^8$.  Damped
variants converge but move in the wrong direction:
\nolinkurl{stage02_convex_pose_velocity_0p01_terminalproj_p0p1_z} leaves
terminal velocity at $7.028\times10^{-6}$ with orders 2.301/1.788, and
\nolinkurl{stage02_convex_pose_velocity_0p01_terminalproj_p0p5_z} leaves
terminal velocity at $1.321\times10^{-5}$ with orders 2.158/1.433.  These
rows set \nolinkurl{projection_used=true}, so they are diagnostic exclusions
rather than candidate acceptance evidence.

\subsection{Why the Next Audit Is an Eight-Row Velocity Compression}

The source-free elimination audit shows that the 24 lower-pair stage rows have
rank 16 after removing the external source; the missing rank is exactly eight
rows.  Mean-velocity and scalar/component blend closures preserve useful order
but leave terminal endpoint velocity at roughly $10^{-5}$.  The final-stage
velocity closure drives terminal velocity to roundoff, but its smooth position
order drops to 3.523.  The problem is therefore not the existence of any
closure row, but the choice of the eight source-free rows that close terminal
velocity while preserving the paper-stage order.

\plotfigure{\figpath/velocity_compression.png}
{Velocity-compression audit.  The all-stage lower-pair velocity rows span the
terminal bridge in the local tangent model, but nonlinear trajectory acceptance
of a universal eight-row compression remains pending.}
{fig:velocity-compression}

The current velocity-compression audit is intentionally local.  It tests 24
candidate/case/step tangent rows and finds that fixed source-free compression
can span the terminal bridge directions, with the best fixed candidate
recording projection residuals around $10^{-15}$.  The strengthened audit now
shows that those fixed compressions collapse to the final-stage selector: the
universal full $8\times24$ selector-relative distance is $1.53\times10^{-15}$
and the largest non-final-stage energy fraction is $1.71\times10^{-15}$.

\plotfigure{\figpath/velocity_compression_nondegenerate.png}
{Non-degenerate velocity-compression probe.  Fixed stage-0/stage-1 injections
with meaningful non-final-stage energy lose the local terminal-bridge span.}
{fig:velocity-compression-nondegenerate}

The follow-up non-degenerate probe makes this failure mode explicit.  It
sweeps 192 candidate/case/step/eta rows by fixing the final-stage coefficient
at one and injecting stage-0 and/or stage-1 lower-pair velocity rows.  There
are 84 rows with meaningful non-final-stage energy, but none of them spans the
local terminal-bridge target; the best meaningful projection residual is
$1.59\times10^{-3}$, far above the $10^{-8}$ local-span threshold.  Only the
eta-zero rows span, and their non-final-stage energy is exactly zero.  Thus the
local compression is not yet a new independent closure formula.

\plotfigure{\figpath/velocity_compression_state_dependent.png}
{State-dependent velocity-compression probe.  Diagonal state-local direction
oracles with meaningful non-final-stage energy still fail to span the local
terminal-bridge target.}
{fig:velocity-compression-state-dependent}

The state-dependent follow-up tests the natural next option: choose diagonal
stage directions from the local tangent state rather than using fixed scalar
stage weights.  It sweeps 72 candidate/case/step/eta rows over a 65-direction
finite grid.  There are 36 rows with meaningful non-final-stage energy, but
none spans the local terminal-bridge target; the best meaningful projection
residual is $9.90\times10^{-3}$.  This rules out diagonal state-dependent
direction weights as the missing independent closure.

\plotfigure{\figpath/velocity_compression_component_mixing.png}
{Non-final component-mixing probe.  Removing the final-stage rows and allowing
full stage-0/stage-1 component mixing still does not span the terminal bridge.}
{fig:velocity-compression-component-mixing}

The full component-mixing follow-up removes the final-stage velocity rows from
the active fit and allows full $8\times8$ or $8\times16$ mixing over stage 0,
stage 1, or both non-final stages.  It tests 36 case/step/candidate rows using
both case-local and universal fits.  None spans the local terminal-bridge
target; the best tangent residual is $9.81\times10^{-1}$.  Some rows are
value-balanced to roundoff, with minimum value residual
$2.19\times10^{-16}$, but value balance alone is not a tangent closure.  The
result motivates the value-level and derivative-aware probes below, rather
than serving as an accepted full-\tfe{} replacement.

\plotfigure{\figpath/velocity_compression_value_level.png}
{Value-level source-free compression probe.  Value fits can balance rows, but
meaningful non-final value-level fits do not span the terminal bridge.}
{fig:velocity-compression-value-level}

The value-level follow-up reverses the diagnostic priority: it fits row values
first, then checks whether the induced Jacobian rows span the terminal bridge.
It tests 72 case/step/candidate rows.  All 72 rows balance values; 18 also
span the local tangent.  However, the spanning value-balanced rows are
final-stage-selector-like: the maximum spanning non-final-stage energy fraction
is $2.12\times10^{-15}$, and the count of meaningful non-final value-balanced
tangent spans is zero.  This rules out a purely value-fitted non-final closure
as the missing independent replacement.  The next target is a meaningful
non-final, derivative-aware value-level compression.

\plotfigure{\figpath/velocity_compression_derivative_aware.png}
{Derivative-aware value-level oracle.  A coefficient-gradient correction gives
meaningful non-final value-balanced tangent spans, but it is still a local
oracle rather than a bounded nonlinear row formula.}
{fig:velocity-compression-derivative-aware}

The derivative-aware follow-up keeps the same value-level non-final row values
and asks whether a local coefficient-gradient correction can close the tangent
gap.  It tests 36 case/step/candidate rows over scalar, component-wise, and
full non-final active velocity families.  All 36 rows balance values and span
the local terminal-bridge tangent with meaningful non-final energy after the
oracle correction.  The best corrected projection residual is
$9.02\times10^{-16}$, while the best uncorrected value-level base residual is
still $9.82\times10^{-1}$ and the maximum required coefficient-gradient
Frobenius norm is $1.95\times10^{17}$.  This is a useful upper bound: it says
the next repair can focus on deriving the missing coefficient-gradient law
rather than searching for another value-only row family.  It is not yet a full
TFE replacement, because the correction has not been converted into a bounded
analytic formula and inserted into a nonlinear trajectory h-sweep.

\plotfigure{\figpath/velocity_compression_bounded_gradient.png}
{Bounded-gradient cap sweep.  The local derivative-aware oracle is capped by a
global Frobenius gradient budget.  Practical caps recover only part of the
local tangent span, while all local spans require cap $10^{18}$.}
{fig:velocity-compression-bounded-gradient}

The bounded-gradient follow-up quantifies how hard the oracle is to realize as
a formula.  It tests 288 cap/candidate/case/step rows.  All 288 rows preserve
the value balance and meaningful non-final row energy, but only 92 rows span at
the tested caps.  A practical cap of $10^{12}$ spans only 12 rows; all 36 local
derivative-aware spans appear only at cap $10^{18}$.  Thus the local tangent
gap is closed only by a very large derivative budget.  The next mathematical
target is still a bounded analytic coefficient-gradient law followed by a
nonlinear h-sweep, not another unconstrained local fit.

\plotfigure{\figpath/velocity_compression_bounded_formula.png}
{Bounded analytic saturation-law probe.  Tanh and rational saturation laws
preserve the value balance, but all local spans require cap $10^{22}$ and the
correction direction is still a target-direction oracle.}
{fig:velocity-compression-bounded-formula}

The bounded-formula follow-up replaces the hard cap with explicit saturation
laws: global tanh, rowwise tanh, and rowwise rational-quadratic scaling.  It
tests 648 law/cap/candidate/case/step rows.  All 648 rows preserve value
balance and meaningful non-final row energy, and 366 rows span at the tested
caps.  For each tested law, complete coverage of the 36 local
derivative-aware spans appears only at cap $10^{22}$.  The best single-row
projection residual is $8.02\times10^{-16}$, attained by the rowwise tanh law
at cap $10^{14}$.  This is still not a full \tfe{} replacement: the saturation
formula bounds the magnitude, but the correction direction remains the local
target-direction oracle and no nonlinear trajectory h-sweep has accepted it.

\plotfigure{\figpath/velocity_compression_target_free_formula.png}
{Target-direction-free bounded formula probe.  The correction direction is
computed from the current closure and active non-final velocity rows, not from
the target Jacobian.}
{fig:velocity-compression-target-free-formula}

The target-direction-free follow-up is the narrowest diagnostic needed after
the saturation-law result.  It tests 2592 law/cap/candidate/case/step rows over
12 direction laws: the positive and negative closure tangent, the mean and
dominant active non-final velocity rows, stage-2 and endpoint velocity
extrapolations from the first two Gauss stages, the non-final stage slope, the
extrapolated-stage2-minus-closure direction, the all-active mean and
anti-aligned rows, and all-active SVD modes with and without closure-row
orthogonalization.  The target bridge is still used to score projection
residuals, but it is not used to construct the correction direction.  All 2592
rows preserve value balance and meaningful
non-final row energy, but none spans the local terminal-bridge
tangent; every tested cap and direction law has zero spanning rows.  The best
projection residual is $9.82\times10^{-1}$, attained by the dominant active
velocity direction at cap $10^{12}$.  This separates two issues that were
previously entangled: the bounded saturation result already showed that a large
target-direction oracle can supply the magnitude, while this target-free probe
shows that the simple source-derived direction laws are not the missing
formula.  As with the preceding probes, this is not yet a full \tfe{}
replacement because it is still a local tangent audit rather than a nonlinear
trajectory h-sweep.

\plotfigure{\figpath/velocity_compression_direction_capacity.png}
{Direction-capacity audit.  All-stage velocity feature spaces can reproduce the
local derivative-aware correction, but non-final active/source feature spaces
cannot.}
{fig:velocity-compression-direction-capacity}

The direction-capacity audit separates feature-space capacity from an actual
target-free formula.  It tests 216 candidate/feature/case/step rows.  The
all-stage velocity feature spaces span 72 rows at roundoff, with best
projection residual $1.09\times10^{-15}$ and best correction residual
$1.51\times10^{-15}$.  Every non-final active/source feature family has zero
spanning rows: the best non-final projection residual is
$9.81\times10^{-1}$ and the best non-final correction residual is
$9.998\times10^{-1}$.  This audit uses the target Jacobian only to define the
local correction being projected; it does not use a target-direction oracle in
the formula path and it is not an accepted full \tfe{} replacement.  It
therefore moves the test from linear feature capacity to nonlinear
second-differential source features.

\plotfigure{\figpath/velocity_compression_nonlinear_capacity.png}
{Nonlinear second-differential capacity audit.  Finite-difference source
features increase the local non-final feature library, but none spans the
local lower-pair correction.}
{fig:velocity-compression-nonlinear-capacity}

The nonlinear-capacity audit tests 216 rows across six finite-difference
source-feature families.  It records zero spanning rows:
the best projection residual is $6.95\times10^{-1}$ and the best correction
residual is $7.01\times10^{-1}$.  The target Jacobian is used only to define
and project the local correction being tested, not to construct a
target-direction oracle formula.  No final-stage velocity selector is accepted,
no nonlinear h-sweep is accepted, and the tested second-differential
non-final source features are ruled out.  This motivates the higher-order and
nonlocal source-feature test below.

\plotfigure{\figpath/velocity_compression_higher_order_capacity.png}
{Higher-order/nonlocal capacity audit.  Third finite-difference source
features and same-case cross-step source deltas still do not span the local
lower-pair correction.}
{fig:velocity-compression-higher-order-capacity}

The higher-order/nonlocal capacity audit tests another 216 rows across six
feature families: third differentials along active non-final velocity rows,
active-plus-third-differential rows, second-plus-third source differentials,
same-case cross-step source-row deltas, active-plus-nonlocal deltas, and a
combined higher-order/nonlocal source feature library.  It again records zero
spanning rows.  The best projection residual is $6.95\times10^{-1}$, and the
best correction residual is $7.01\times10^{-1}$.  This closes the narrow next
step from the previous audit: the missing coefficient-gradient law is not
present in these tested third-differential or simple cross-step nonlocal
source features.  The next repair target is therefore a different
history-dependent non-final source law or a revised weak-row formula, followed
by the nonlinear h-sweep.

\plotfigure{\figpath/velocity_compression_history_capacity.png}
{History-transport capacity audit.  One-step transported source-row and
non-final velocity-row deltas still do not span the local lower-pair
correction.}
{fig:velocity-compression-history-capacity}

The history-transport capacity audit tests another 216 rows across six
feature families: transported source-row deltas, active-plus-source transport
deltas, transported non-final velocity-row deltas, active-plus-velocity
transport deltas, combined source/velocity transport deltas, and the same
combined transport library with active non-final velocity rows included.  It
again records zero spanning rows.  The best projection residual is
$7.13\times10^{-1}$, and the best correction residual is
$7.12\times10^{-1}$.  Thus a one-step trajectory-history feature space is not
enough.

\plotfigure{\figpath/velocity_compression_multi_step_history_capacity.png}
{Two-step history-transport capacity audit.  Transported source-row and
non-final velocity-row deltas from two future states still do not span the
local lower-pair correction.}
{fig:velocity-compression-multi-step-history-capacity}

The two-step history-transport capacity audit advances the terminal output
twice and stacks the transported source-row and non-final velocity-row deltas
from both future states.  It again tests 216 rows across six feature families
and records zero spans.  The best projection residual is
$6.95\times10^{-1}$, and the best correction residual is
$7.01\times10^{-1}$.  The remaining target is therefore no longer simply
``try history'': the tested one-step and two-step transported feature spaces
are insufficient.

\plotfigure{\figpath/velocity_compression_recurrent_history_capacity.png}
{Recurrent history-capacity audit.  Two-step curvature and bilinear
source/velocity history features still do not span the local lower-pair
correction.}
{fig:velocity-compression-recurrent-history-capacity}

The recurrent history-capacity audit then tests 216 rows over two-step
source/velocity curvature and pairwise bilinear source/velocity history
features.  It records zero spans; the best projection and correction residuals
are $6.952\times10^{-1}$ and $7.012\times10^{-1}$.  The next repair therefore
needs a revised weak-row formula or richer nonlinear recurrent history source
law, followed by the nonlinear h-sweep.

\plotfigure{\figpath/velocity_compression_weak_row_structure_capacity.png}
{Weak-row structure-capacity audit.  Source/active cross-Gram, Hadamard,
cross-Hadamard, and shifted-commutator feature rows still do not span the local
lower-pair correction.}
{fig:velocity-compression-weak-row-structure-capacity}

The weak-row structure-capacity audit then tests 216 rows over six algebraic
source-active feature families.  These features are target-free in construction
and use the target Jacobian only for capacity scoring.  They again record zero
spans; the best projection and correction residuals are
$6.952\times10^{-1}$ and $7.012\times10^{-1}$.  This is a narrower negative
result than the history tests: the missing row is not recovered by the tested
cross-Gram, Hadamard, cross-Hadamard, or shifted-commutator structure rows.
The next repair therefore needs a revised analytical weak-row formula or a
richer nonlinear recurrent history source law.

\plotfigure{\figpath/velocity_compression_row_space_compression.png}
{Row-space compression audit.  Target-free SVD, row-norm, and stage-balanced
frozen eight-row compressions of the all-stage/non-final velocity row space do
not provide a value-balanced tangent closure.}
{fig:velocity-compression-row-space-compression}

The row-space compression audit removes the target Jacobian and target
direction from the formula construction.  It builds frozen case-local
compression matrices from the velocity row space only, then scores whether the
resulting eight rows span the terminal lower-pair bridge.  Across 36 rows
covering six laws, two cases, and three step sizes, it records zero spans and
24 value-balanced rows.  The best projection residual is
$7.378\times10^{-1}$ for \artifact{all\_stage\_svd\_top8}; the best value
residual is $3.68\times10^{-12}$ for the same law, but that value balance does
not close the tangent gap.  The best non-final residual is
$9.815\times10^{-1}$.  Because the coefficients are frozen state-dependent
quantities and coefficient derivatives are not included, this is not a
nonlinear full-\tfe{} residual.  It rules out this target-free row-space
compression family and leaves the same repair target: a revised analytical
weak-row formula or richer nonlinear recurrent history source law.

\plotfigure{\figpath/order_closure_blend.png}
{Order/closure blend trajectory audit.  The beta path connects the
order-preserving mean-acceleration closure to the terminal-closing final-stage
velocity closure.  It confirms that the two properties remain split.}
{fig:order-closure-blend}

That targeted promotion has now been run as an 18-row diagnostic over two
cases, three step sizes, and beta values $\{0,0.5,1\}$.  The beta-zero and
beta-half smooth branches preserve position order above five, but leave raw
terminal endpoint velocity around $4\times10^{-6}$.  The beta-one branch closes
terminal velocity to $4.01\times10^{-16}$, but has smooth position order 3.523.
Thus the intersection flag is false: no tested beta simultaneously satisfies
the terminal-velocity and order gates.  The evidence narrows the next repair to
a genuinely new eight-row source-free closure formula rather than another
scalar blend on the same endpoints.

A newer recurrent/terminal gamma-blend smoke promotes the same tradeoff to a
bounded trajectory h-sweep without invoking the full generator.  The latest
default run used smooth and sharp cases, $\gamma\in\{0.5,1.0\}$,
$h\in\{0.04,0.02\}$, and an $h=0.01$ reference.  It converged all 24 requested
steps with rank 132 and maximum residual $4.77\times10^{-12}$.  The
terminal-closing branch is still $\gamma=1.0$, with raw terminal endpoint
velocities $3.05\times10^{-16}$ and $2.94\times10^{-16}$ in the smooth and
sharp cases, but its smooth minimum order is only 4.397 and the sharp velocity
order remains $-0.428$.  The best smooth-order branch, $\gamma=0.5$, reaches
minimum order 4.413 but leaves terminal velocity up to $8.11\times10^{-6}$ on
the sharp case.  This targeted result is consistent with the larger
order/closure audit: the properties remain split, so the full-\tfe{} gate stays
open.

A denser bounded scalar-gamma screen over
$\gamma\in\{0,0.25,0.5,0.75,1\}$ converged all 60 requested short-trajectory
steps with the same rank-132 residual structure.  It did not change the
conclusion.  The best smooth-order branch moved only to $\gamma=0.75$, with
minimum order 4.414, while terminal closure remained only on the
$\gamma=1.0$ branch and the maximum terminal velocity over the grid stayed at
$8.23\times10^{-6}$.  The scalar interpolation family is therefore exhausted
as a plausible repair; the next candidate must change the eight lower-pair
closure rows themselves.

A near-terminal scalar-policy screen over
$\gamma\in\{0.9,0.99,0.999,1\}$ also fails.  It converges all 48 requested
short-trajectory steps, and its best smooth-order branch is $\gamma=0.99$ with
minimum order 4.415.  However, that branch leaves sharp terminal velocity at
$3.08\times10^{-6}$, while even $\gamma=0.999$ leaves
$4.68\times10^{-7}$.  Only $\gamma=1.0$ closes terminal velocity to roundoff,
and that is exactly the order-limited final-stage endpoint.  This closes the
near-terminal scalar-policy loophole.

A component-wise one-step gamma sensitivity check then releases one of the
eight terminal directions at a time to $\gamma=0.999$.  The run covers smooth
and sharp cases with 10 gamma vectors, for 20 one-step rows, and completes in
31.6 seconds.  It preserves rank 132 with maximum residual
$1.77\times10^{-12}$.  Five nontrivial rows close terminal velocity at the
$10^{-12}$ tolerance; the best is
\texttt{component\_2\_release\_0p999} at $3.46\times10^{-17}$, while
\texttt{component\_7\_release\_0p999} is the worst component release at
$1.23\times10^{-11}$.  This is useful sensitivity evidence, but it is still a
one-step screen rather than a trajectory h-sweep, so the independent
full-\tfe{} replacement remains open.

The same component idea was then promoted to a bounded trajectory h-sweep over
four vector labels:
\texttt{all\_ones}, \texttt{all\_0p999},
\texttt{component\_2\_release\_0p999}, and
\texttt{component\_7\_release\_0p999}.  The run covers smooth and sharp cases,
$h\in\{0.04,0.02\}$, and an $h=0.01$ reference, converging all 48 short
trajectory steps in 98.4 seconds with rank 132 and maximum residual
$4.77\times10^{-12}$.  The best terminal vector is
\texttt{component\_2\_release\_0p999}, with terminal velocity
$2.70\times10^{-16}$, but its smooth minimum order remains about 4.397.  The
best smooth-order vector is \texttt{all\_0p999}, with smooth minimum order
4.401, but it leaves terminal velocity $4.68\times10^{-7}$.  Thus the
trajectory-level component screen still has no order/terminal intersection.

A final lightweight nonlinear-source smoke sets the component-wise blend from
the previous step's history source, rather than from a prescribed scalar or
fixed vector.  The default smooth-only
\texttt{history\_abs\_inf} run uses releases $\{0.001,0.01\}$, the same
$h\in\{0.04,0.02\}$ sweep, and an $h=0.01$ reference.  It converges all 12
short-trajectory steps in 38.9 seconds with rank 132 and maximum residual
$2.59\times10^{-12}$.  The best branch has terminal velocity
$1.45\times10^{-7}$ and smooth minimum order 4.402.  A signed-policy follow-up
with \texttt{history\_signed\_positive} and
\texttt{history\_signed\_negative} converges in 36.9 seconds; its best
terminal velocity is still $1.45\times10^{-7}$ and its best smooth minimum
order is 4.415.  This excludes the simple explicit recurrent history-gamma
law as the missing full-\tfe{} replacement.  The next target is therefore a
revised analytical weak-row formula or a richer nonlinear history source law
that changes the eight closure rows themselves.

The next diagnostic looks at that weak-row formula locally.  A one-step
recurrent weak-row differential audit evaluates the current source-free
recurrent weak closure in the smooth case at $h=0.04$, using both the
zero-initial and post-one-step history modes.  It completes in 26.0 seconds,
keeps rank 132, and has maximum residual $2.48\times10^{-12}$.  However, the
current recurrent weak row does not span the terminal-bridge tangent: the
projection residual remains between 0.899 and 0.974 and the independent target
rank remains eight.  The final-stage terminal row, used only as a scoring
comparison, projects at $1.46\times10^{-15}$.  The missing direction is
stage-2-local and dominated by translational velocity variables, while the
relative recurrent-vs-final closure Jacobian gap reaches 6.872.  This is the
most concrete current blocker: the next formula must supply that stage-2
velocity coefficient-gradient contribution without substituting the
final-stage terminal row.

A follow-up source/history coefficient-gradient probe tests exactly that
direction without running a full trajectory solve.  It modulates the stage-2
lower-pair velocity row by curvature-plus-history features with gains
$10^8$ and $10^{12}$.  The smooth local run covers both history modes and 8
candidate rows, completes in 35.3 seconds, keeps rank 132, and has maximum
residual $2.48\times10^{-12}$.  None of the candidates spans the
terminal-bridge tangent: the best projection residual remains 0.899 and the
independent target rank remains eight.  The largest closure-value perturbation
from the recurrent weak row is $1.25\times10^{-9}$, showing that this simple
stage-2 coefficient-gradient correction perturbs the row value without adding
the missing tangent.  The next formula must therefore be richer than a scalar
source/history modulation of the stage-2 velocity row.

A matrix-mixing follow-up makes that statement stronger while still avoiding a
full trajectory solve.  It applies target-free row-broadcast, column-broadcast,
symmetric-broadcast, diagonal-plus-row-broadcast, shifted-column-broadcast, and
bilinear stage-2 velocity coefficient matrices to the same stage-2 lower-pair
velocity row, with gains $10^8$ and $10^{12}$.  The smooth local run covers
both history modes and 32 candidate rows, completes in 68.1 seconds, keeps rank
132, and has maximum residual
$2.48\times10^{-12}$.  None of these matrix-mixed rows spans the
terminal-bridge tangent either: the best law is
\texttt{diagonal\_plus\_row\_broadcast\_feature} at gain $10^{12}$, with
projection residual 0.898873 and independent target rank eight.  Its largest
closure-value perturbation from the recurrent weak row is
$4.00\times10^{-9}$, so row/column mixing changes the value but still does not
supply the missing stage-2 tangent.  The remaining full-\tfe{} repair target is
therefore a genuinely different analytical weak-row formula, not a matrix
premultiplication of the current stage-2 velocity row.

An active-velocity extension of the same matrix audit then uses the stage-2
velocity row itself, the stage-2 minus stage-1 velocity row, and the velocity
curvature as coefficient features.  This run again covers both history modes,
16 candidate rows, and gains $10^8$ and $10^{12}$; it completes in 46.6
seconds with rank 132 and maximum residual $2.48\times10^{-12}$.  The best
law, \nolinkurl{stage2_velocity_column_broadcast_feature} at gain
$10^{12}$, improves the local projection residual to 0.585746, but still has
independent target rank eight and no terminal-bridge span.  Its largest
closure-value perturbation is $1.46\times10^{-6}$.  Thus active stage-2
velocity dependence carries more of the missing tangent than source/history
features, but it still cannot be accepted as a full-\tfe{} replacement or
promoted to a trajectory h-sweep without a different weak-row formula.

A final small active-law grid widens that check without adding any new
trajectory machinery.  It tests symmetric, diagonal-plus-row, and shifted
column broadcasts of the stage-2 velocity feature, plus column broadcasts of
stage-2-minus-stage-1 and stage-0/1 mean velocity.  The grid covers 20 local
rows and completes in 51.2 seconds.  The best row is
\nolinkurl{stage2_velocity_shifted_column_broadcast_feature} at gain
$10^{12}$, improving the projection residual only to 0.576548.  The
independent target rank remains eight and no row spans the terminal bridge.
This leaves the active-velocity family useful as localization evidence, not as
an acceptably derived eight-row lower-pair closure.

The best active-law row was then rerun alone with a missing-direction
decomposition, again without invoking the full artifact pipeline.  This
targeted check covers 2 local rows in 27.8 seconds.  The zero-initial row has
projection residual 0.909693; its missing direction has rank eight, dominant
variable family \nolinkurl{translation_velocity_v} at fraction 0.487, and
stage-2 fraction 0.965.  The post-one-step row is the best row, with
projection residual 0.576548; its missing direction again has rank eight, but
the dominant variable family shifts to \nolinkurl{angular_velocity_w} at
fraction 0.503, with stage fractions 0.0185, 0.0185, and 0.963 for stages 0,
1, and 2.  This makes the next repair target narrower: it should derive a
stage-2 lower-pair weak row whose coefficient derivative couples the velocity
closure to angular velocity, rather than continuing to widen generic
matrix-premultiplication families.

A stage-2 angular-velocity mask follow-up tests the simplest version of that
localization: keep only the angular components of the stage-2 lower-pair
velocity feature.  It covers 8 local rows in 35.3 seconds with four matrix
laws and gain $10^{12}$.  The best law,
\nolinkurl{stage2_angular_velocity_shifted_column_broadcast_feature}, reaches
projection residual 0.752997, worse than the unmasked active-law residual
0.576548.  The independent target rank remains eight and no row spans the
terminal bridge.  The best missing direction shifts back to
\nolinkurl{translation_velocity_v} at fraction 0.660, with 97.8\% of the
energy still in stage 2.  Thus the missing formula is not a simple angular-only
mask; it likely needs coupled stage-2 translational/angular weak-row
structure.

A direct translation/angular cross-coupling follow-up then tests the simplest
outer-cross matrix versions of that idea.  It covers 8 local rows in 35.3
seconds with four laws and gain $10^{12}$.  The best law,
\nolinkurl{stage2_velocity_angular_to_translation_cross_feature}, reaches
projection residual 0.898812, again worse than the unmasked active-law
residual 0.576548.  The independent target rank remains eight and no row spans
the terminal bridge.  Its remaining direction is dominated by
\nolinkurl{translation_velocity_v} at fraction 0.556, with essentially all
energy in stage 2.  Thus simple outer-cross coefficient matrices are also not
the missing weak row.

A target-free endpoint-pose velocity-predictor follow-up moves closer to the
missing row family.  It evaluates lower-pair velocity constraints at the paper
terminal pose using several generalized-velocity predictors from the stage
data.  The best of 10 local rows uses the positive part of the Gauss-node
terminal Lagrange velocity predictor.
It reduces the projection residual to 0.117782, but still leaves independent
target rank eight and no terminal-bridge span.  The remaining direction is
dominated by \nolinkurl{translation_velocity_v} at fraction 0.514, with 96.2\%
of its energy in stage 0.  This is the sharpest current local repair probe,
but it remains diagnostic.

A Gauss endpoint-pose version of the same predictor family is negative.  The
best of 7 local rows is
\nolinkurl{gauss_endpoint_pose_positive_lagrange_z}; it reaches only residual
0.208599, leaves independent target rank eight and no span, and changes the
dominant missing family to \nolinkurl{lie_position_u} at fraction 0.651.  Since
this is worse than the paper endpoint-pose residual 0.117782, it is a ruled-out
endpoint-pose alternative rather than a new repair direction.
The remaining 11 Gauss endpoint-pose quarter, stage1-half, convex, and mean
predictors are also negative.  Their best row,
\nolinkurl{gauss_endpoint_pose_stage02_convex_0p00_z}, reaches residual
0.172659, leaves independent target rank six and no span, and keeps the
remaining direction dominated by \nolinkurl{lie_position_u} at fraction
0.952340.

A near-terminal stage-0/stage-2 convex predictor screen clarifies the remaining
boundary.  The stage-2-only endpoint-pose row spans the terminal bridge at
roundoff, but this is exactly the terminal-stage bridge row already known to be
order-limited in nonlinear h-sweeps.  The best nonterminal row uses a 5\%
stage-0, 95\% stage-2 velocity predictor; it lowers the local residual to
0.049317 at $h=0.04$, but still leaves independent target rank eight and no
nonterminal span.  Focused reruns at $h=0.02$ and $h=0.01$ give residuals
0.051890 and 0.052472.  The terminal-equivalent row still spans at roundoff,
but the nonterminal row does not converge to a span as $h$ is reduced, so the
remaining target is an analytical nonterminal weak row rather than a
continuation toward the terminal bridge.

The synchronized pose/velocity version of the same test evaluates
$C_v(q_\alpha,z_\alpha)$ at matching stage-0/stage-2 convex points.  Its best
row, \nolinkurl{stage02_convex_pose_velocity_0p01_z}, reduces the local
residual to 0.009512 at $h=0.04$.  However, it still leaves rank eight
unspanned, and h-scaling gives 0.009978 at $h=0.02$ and 0.010085 at
$h=0.01$.  The residual is nearly proportional to the offset from the
terminal row: 0p02 gives 0.019205, 0p05 gives 0.049390, and 0p10 gives
0.103464.  This is useful localization evidence but still a near-terminal
limit, not an independent full-\tfe{} replacement.

A terminal-limit extrapolation check then combines two nonterminal synchronized
rows.  The 0p01/0p02 extrapolated law reduces the local residual to
$4.699\times10^{-6}$ at $h=0.04$, and the wider 0p01/0p05 extrapolation gives
$1.175\times10^{-5}$.  Under h-scaling the 0p01/0p02 residual decreases to
$2.335\times10^{-6}$ and $1.168\times10^{-6}$ at $h=0.02$ and $h=0.01$, but
it still does not span the terminal bridge: the independent target ranks are
8, 8, and 7.  We therefore mark these rows as terminal-bridge-equivalent
extrapolations, not as independent full-\tfe{} replacements.

Adding the centered bridge's mean-acceleration term directly does not fix the
remaining tangent.  The signed acceleration-bridge variants worsen the
projection residual from $4.699\times10^{-6}$ to at least 0.076159 and move
the dominant missing direction to \nolinkurl{lie_position_u}.  This excludes
the direct mean-acceleration bridge term as the missing nonterminal correction.

A bilinear active-velocity outer-product follow-up then checks whether simple
component-pair coupling between the coefficient feature and the stage-2
velocity row supplies the missing tangent.  It covers 24 local rows in 57.3
seconds using feature/stage-2-velocity outer, reverse outer, symmetric outer,
and related active-feature variants.  The best row is
\nolinkurl{stage2_velocity_feature_outer_stage2_velocity} at gain $10^{12}$,
with projection residual 0.898720.  The independent target rank remains eight
and no row spans the terminal bridge.  This rules out the tested bilinear
outer-product coupling family; it does not improve on the previous
shifted-column active-law localization residual 0.576548.

A componentwise active-velocity diagonal follow-up then tests whether the
missing stage-2 tangent is row-local rather than an outer product.  It covers
24 local rows in 56.7 seconds using Hadamard diagonal, shifted diagonal, and
diagonal-plus-row/column variants.  The best row is
\nolinkurl{stage2_velocity_diagonal_plus_hadamard_row_feature_stage2_velocity}
at gain $10^{12}$, with projection residual 0.898530.  The independent target
rank remains eight and no row spans the terminal bridge.  This rules out the
tested componentwise and shifted-diagonal active-velocity couplings as the
missing eight-row weak closure.

Two acceleration-correction follow-ups test whether the synchronized
near-terminal row only needs a simple Taylor term.  The first adds the centered
bridge mean-acceleration closure directly to
\nolinkurl{stage02_convex_pose_velocity_0p01_z} and its 0p02 analogue with
signed coefficients $\pm0.1$ and $\pm1$.  The second applies the same signed
coefficient as a generalized-velocity shift,
$z_\alpha+\gamma h\dot z_\alpha$, before evaluating $C_v(q_\alpha,z_\alpha)$.
Each grid covers 18 local rows, keeps rank 132, and records zero nonterminal
spans.  In both grids the best row remains the uncorrected
\nolinkurl{stage02_convex_pose_velocity_0p01_z} with residual 0.009512; the
best acceleration-corrected residual is only 0.012376, while unit corrections
are no better than 0.077874.  Thus a direct acceleration closure or
generalized-velocity acceleration predictor is not the missing nonterminal
full-\tfe{} weak row.

A further geometry-oriented test replaces or perturbs $z_\alpha$ toward the
stage-0/stage-2 pose slope $(q_2-q_0)/(h(c_2-c_0))$.  This is meant to check
whether the missing row is just a kinematic velocity estimate from the paper
stage positions.  It covers 26 local rows, again with no nonterminal span.  The
best row is still the uncorrected
\nolinkurl{stage02_convex_pose_velocity_0p01_z} at residual 0.009512; the best
pose-slope row,
\nolinkurl{stage02_convex_pose_velocity_0p02_poseslope_m0p1_z}, has residual
0.878519 and leaves a stage-2 dominated translation-velocity direction.  The
tested pose-slope velocity predictor is therefore also not the missing
full-\tfe{} weak row.

Finally, a source-to-velocity lift test maps lower-pair source rows through
the current endpoint velocity matrix $C_{v,z}(q_\alpha)$ into a minimum-norm
generalized-velocity correction.  This is a richer target-free source law than
the scalar Taylor corrections: it uses the mean acceleration closure, stage
source estimates, source curvature, and the history/source delta.  The full
grid covers 34 local rows in 76.4 seconds and again has zero nonterminal spans.
The best row overall is still the uncorrected
\nolinkurl{stage02_convex_pose_velocity_0p01_z} at residual 0.009512.  A
focused rerun gives the best source-lift residual 0.609803 for
\nolinkurl{stage02_convex_pose_velocity_0p01_sourcelift_historydelta_p1_z}
and the equivalent source0 branch.  This excludes the tested current-pose
$C_v$ source-to-velocity lift as the independent weak row.

A final near-terminal slope probe tests the finite-difference slope of the
same synchronized pose/velocity family rather than its terminal-limit
intercept.  It covers four local rows in 25.8 seconds, keeps rank 132, and
again has no span.  The best law,
\nolinkurl{stage02_convex_pose_velocity_slope_0p01_0p05_z}, reaches only
projection residual 0.677034 with independent target rank eight and dominant
remaining family \nolinkurl{translation_velocity_v}.  Thus the first
nonterminal slope of the near-terminal convex row is also not the missing
coefficient-gradient closure.

A near-terminal curvature probe then tests the second finite-difference
curvature of the same synchronized pose/velocity family.  It covers four local
rows in 26.7 seconds, keeps rank 132, and again has no span.  The best law,
\nolinkurl{stage02_convex_pose_velocity_curvature_0p01_0p05_0p10_z}, reaches
only projection residual 0.871228 with independent target rank eight and
dominant remaining family \nolinkurl{translation_velocity_v}.  Thus the local
second derivative of the near-terminal convex row is also not the missing
coefficient-gradient closure.

A source/history coefficient-feature matrix probe then tests recurrent
features such as history delta, source0/source2, source02 delta, and source
curvature as shifted-column coefficient matrices multiplying the stage-2
lower-pair velocity row.  It covers 16 local rows in 44.8 seconds, keeps rank
132, and again has no span.  The best law,
\nolinkurl{source_curvature_shifted_column_broadcast_feature}, reaches only
projection residual 0.898873 with independent target rank eight and dominant
remaining family \nolinkurl{translation_velocity_v}.  Thus the tested
source/history coefficient features are also not the missing
coefficient-gradient closure.

A nonlinear recurrent curvature/history matrix-feature extension then adds
source-curvature minus history-delta, source/history Hadamard, unit-Hadamard,
and normalized curvature-plus-history feature prefixes.  It covers 16 local
rows in 45.5 seconds, keeps rank 132, and again has no span.  The best law,
\nolinkurl{source_curvature_minus_history_delta_diagonal_plus_row_broadcast_feature},
reaches only projection residual 0.898865 with independent target rank eight
and a stage-2 dominated \nolinkurl{translation_velocity_v} gap, with dominant
fraction 0.556051 and stage-2 fraction 0.999996.  Thus the tested nonlinear
recurrent curvature/history matrix features are also excluded.

A terminal-limit source-lift matrix-difference probe then tests whether the
best near-terminal extrapolated intercept can be repaired by adding a
stage-local $C_v$ source-to-velocity lift before extrapolation.  It covers
eight local rows in 47.5 seconds, keeps rank 132, and again has no span.  The
best law is the \texttt{curvature\_m0p1} source-lift variant; it reaches
projection residual $5.298\times10^{-6}$ with independent target rank eight
and dominant remaining family \nolinkurl{translation_acceleration_a}.
This is close to, but not better than, the uncorrected terminal-limit residual
$4.699\times10^{-6}$, so the tested $C_v$ source-lift correction is also not
the missing coefficient-gradient closure.
A curvature source-lift scale refinement then tests the same terminal-limit
row with signed curvature coefficients from $0.01$ to $0.5$.  It covers eleven
local rows in 52.9 seconds and again has no span; the best
\texttt{curvature\_p0p01} law still reaches only
$5.298\times10^{-6}$, with independent target rank eight, dominant
\nolinkurl{translation_acceleration_a} fraction 0.407672, and stage fractions
0.549481/0.005742/0.444777.  The local source-curvature norm is
$1.076\times10^{-14}$, so this is not merely a missed source-lift scale.
A mean-acceleration bridge scale refinement then tests the same terminal-limit
row with signed bridge coefficients from $0.01$ to $0.5$.  It covers eleven
local rows in 31.4 seconds and again has no span.  The best nonterminal law,
\texttt{accel\_p0p01}, reaches only 0.000867721 with independent target rank
eight and dominant \nolinkurl{lie_position_u} fraction 0.706665, while the
overall best row remains the terminal-bridge-equivalent uncorrected
extrapolation.  Thus this is not merely a missed small acceleration-bridge
coefficient.
A generalized-acceleration Taylor-shift scale refinement then uses separate
$h\dot z$ shifts at the left and right convex points before the same
extrapolation.  It covers eleven local rows in 30.2 seconds and again has no
span.  The best nonterminal law, \texttt{genaccel\_p0p01}, reaches only
0.000866899 with independent target rank eight and dominant
\nolinkurl{lie_position_u} fraction 0.701135.  Thus the missing row is not the
tested stage-local generalized-acceleration Taylor shift.
A pose-acceleration Taylor-shift scale refinement then moves the evaluated
pose rather than only the velocity, using separate $h^2\dot z$ shifts at the
same left and right convex points.  It covers eleven local rows in 29.9
seconds and again has no span.  The best nonterminal law,
\texttt{poseaccel\_p0p01}, reaches 4.1835e-05 with independent target rank
eight, dominant \nolinkurl{lie_position_u} fraction 0.711570, and stage-2
fraction 0.771337.  This pose-directed perturbation is closer than the
velocity-only generalized-acceleration shift, but it still does not provide an
independent full-\tfe{} stage row.
A tiny pose-acceleration scale refinement then lowers the signed coefficients
to $0.0005$, $0.001$, $0.002$, and $0.005$, while retaining the previous
positive $0.01$ row.  It covers ten local rows in 30.0 seconds and still has no
span.  The best nonterminal law,
\texttt{poseaccel\_p0p0005}, reaches 5.5898e-06 with independent target rank
eight and dominant \nolinkurl{translation_acceleration_a} fraction 0.350240.
Thus small pose shifts only approach the terminal-equivalent row and do not
produce an independent full-\tfe{} closure row.
A pose+velocity acceleration Taylor refinement then combines the small
$h^2\dot z$ pose shift with signed small $h\dot z$ velocity shifts.  It covers
eleven local rows in 30.7 seconds and again has no span.  The best nonterminal
law, \texttt{posevelaccel\_p0p0005\_m0p0001}, reaches 1.0173e-05 with
independent target rank eight and dominant \nolinkurl{lie_position_u} fraction
0.544660.  Thus this simple combined Taylor correction is also not the missing
independent full-\tfe{} closure row.
A component-split pose-acceleration refinement then separates the same
$h^2\dot z$ pose shift into translation-only and angular/Lie-only branches.
The two JSON-only screens cover 26 local rows in 32.3+32.3 seconds, keep rank
132, and again have no span.  Translation-only shifts through coefficient
0.1 are locally indistinguishable from the uncorrected terminal-limit row; the
best nonterminal law, \texttt{transposeaccel\_m0p01}, remains at 5.298e-06 with
dominant \nolinkurl{translation_acceleration_a} fraction 0.407672.  The
angular-only branch reproduces the earlier full pose-acceleration behavior:
\texttt{angposeaccel\_p0p002} reaches 9.6511e-06, and
\texttt{angposeaccel\_p0p01} reaches 4.1835e-05.  Thus the observed
pose-acceleration effect is angular/Lie-pose dominated; a translation-only
pose shift is not the missing independent weak row.

A component-split velocity-acceleration Taylor refinement then performs the
analogous split on the $h\dot z$ generalized-velocity shift.  The first
JSON-only screen covers 13 local rows in 31.8 seconds, keeps rank 132, and
finds no span.  Its best nonterminal row,
\texttt{transvelaccel\_m0p01}, reaches projection residual
$2.138\times10^{-4}$ with independent target rank eight and dominant
\nolinkurl{translation_acceleration_a} fraction 0.886603.  The best angular
velocity branch, \texttt{angvelaccel\_p0p01}, reaches only 0.000864560 and
moves the dominant missing family to \nolinkurl{lie_position_u}.  A tiny
translational scale follow-up covers 11 local rows in 31.4 seconds; its best
law, \texttt{transvelaccel\_m0p0005}, reaches
$9.8175\times10^{-6}$, still worse than the uncorrected terminal-limit
residual $5.298\times10^{-6}$.  Thus the velocity-acceleration split only
approaches the terminal-equivalent limit and does not provide an independent
full-\tfe{} closure row.

A component-mixed pose/velocity Taylor refinement then combines the two most
promising split directions: angular/Lie pose shift with translational velocity
shift, plus the complementary translation-pose/angular-velocity branch.  It
covers 13 local rows in 32.0 seconds, keeps rank 132, and again finds no span.
The best nonterminal row is the $p0.0005/m0.0005$ angular-pose/translational-
velocity row.  It reaches only $9.987\times10^{-6}$, with independent target
rank eight and a dominant translation-acceleration fraction of 0.774582.  The
complementary translation-pose/angular-velocity branch reaches only
0.000864560.  This is still worse than the uncorrected terminal-limit residual
$5.298\times10^{-6}$, so the tested component-mixed Taylor cross term is also
not the missing independent full-\tfe{} closure row.

A stage-2-fixed/delta acceleration velocity-shift refinement then tests whether
the dominant translation-acceleration residual is an artifact of using the
interpolated acceleration source.  It covers 13 local rows in 31.8 seconds,
keeps rank 132, and again finds no span.  The best nonterminal row is
\texttt{transvelaccelstage2\_m0p0005}, with residual
$9.8175\times10^{-6}$ and translation-acceleration fraction 0.802349.  The best
delta20 row reaches only $1.0119\times10^{-5}$.  Thus stage-2-fixed and
delta20 acceleration sources reproduce the earlier tiny-shift limit rather
than providing the independent full-\tfe{} closure row.

A nonfinal terminal velocity/source predictor refinement then tries to predict
the terminal lower-pair velocity closure from only stage-0/stage-1 velocity rows
and source estimates.  It covers nine local rows in 30.9 seconds, keeps rank
132, and finds no span.  The best law, \texttt{euler1}, reaches only 0.923466,
with dominant \nolinkurl{translation_velocity_v} fraction 0.526951 and stage-2
fraction 0.959959.  Thus a plain nonfinal velocity/source integral predictor
does not supply the missing stage-2 terminal-bridge Jacobian direction.
The corresponding bounded trajectory smoke confirms this at nonlinear-solve
level: the stage-0/stage-1 family converges 24/24 rows with rank 132, but its
best terminal velocity is only $1.770\times10^{-7}$, and the order-friendlier
\texttt{euler1}/\texttt{ab01}/\texttt{source01linear1} rows remain at
$3.584\times10^{-7}$.  A stage-0/1/2 extension using the third Gauss stage also
converges all 18 rows with \nolinkurl{projection_used=false}, but
\nolinkurl{terminal_velocity_closed=false}: the closest rows reach
$2.559\times10^{-11}$ at $h=0.04$ and then jump to $1.328\times10^{-7}$ at
$h=0.02$, while the stable \texttt{linear012} row is only
$1.134\times10^{-7}$.

A stage-2 source-to-velocity transport refinement applies the source/history
projection directly at the stage-2 pose/velocity closure.  The first screen
covers 13 local rows in 53.4 seconds and reaches 0.001251 with independent
target rank six.  A scale refinement over the matrix-difference history-delta
family covers 11 rows in 57.7 seconds and improves the best residual to
0.0001251 at \texttt{historydelta\_m0p01}.  An ultra-fine scale follow-up
covers nine rows in 49.6 seconds and pushes the best nonterminal residual to
$1.251\times10^{-6}$ at \texttt{historydelta\_m0p0001}.  This is the strongest
nonterminal transport clue so far, but it still has no span and leaves a
rank-six \nolinkurl{lie_position_u} missing direction.  A combined
source-transport/angular-pose follow-up covers eight rows in 48.7 seconds; the
source-only row remains best, and the best combined row reaches only
$1.296\times10^{-6}$ while raising the independent target rank to eight and
still finding no span.
Promoting those two source-transport rows into the bounded nonlinear trajectory
smoke does not change the acceptance boundary.  The three-law comparison runs
in 74.1 seconds, keeps rank 132, uses no projection, and writes no artifacts.
The source-only transported row
\nolinkurl{stage02_convex_pose_velocity_0p00_sourceliftmatrixdiff_historydelta_m0p0001_z}
keeps the two-point trajectory orders at 6.588/4.508, but its max terminal
velocity is $8.020\times10^{-12}$, above the $10^{-12}$ closure tolerance.
Adding the small angular-pose shift worsens the terminal velocity to
$2.102\times10^{-7}$ and lowers the velocity order to 2.820.  Thus the
source-to-velocity transport clue is now a trajectory-level negative as well
as a local no-span result.
A three-$h$ refinement of the source-only transported row confirms that this
is not a two-point artifact: it converges all seven short-sweep steps in 46.4
seconds, keeps rank 132 and no projection, but its maximum terminal velocity is
still $8.020\times10^{-12}$ even though the finest-$h$ value is
$5.904\times10^{-13}$.  Its position/velocity orders are again only
4.142/2.305, matching the bare endpoint-row order loss rather than recovering
the full-\tfe{} acceptance target.
Shrinking the same source-transport coefficient by another factor of ten gives
a useful limiting negative: the
\nolinkurl{stage02_convex_pose_velocity_0p00_sourceliftmatrixdiff_historydelta_m0p00001_z}
row closes the terminal tolerance with max/finest terminal velocity
$8.021\times10^{-13}$/$5.910\times10^{-14}$, but the smooth orders remain
4.142/2.305.  Coefficient shrinking can therefore recover terminal closure in
this family, but only by approaching the same order-limited endpoint-row
behavior.

\section{Next Full-\tfe{} Repair Target}

The remaining full-\tfe{} work is now narrow enough to state as a concrete
repair target.  It is not another four-example validation run, and it is not a
request to rerun the complete historical audit harness.  The missing object is
an eight-row, source-free lower-pair closure formula that can replace the
terminal bridge while preserving the trajectory order of the paper-derived
stage residual.

\begin{table}[t]
\centering
\caption{Acceptance test for a future independent full-\tfe{} replacement.}
\label{tab:full-tfe-acceptance-test}
\begin{tabular}{P{0.27\linewidth}P{0.61\linewidth}}
\toprule
Requirement & Acceptance evidence \\
\midrule
Stage residual replacement &
The nonlinear solve uses the paper-derived 132 stage rows, including the new
eight lower-pair closure rows, rather than appending an endpoint correction to
a Gauss trajectory. \\
Source-free construction &
The lower-pair closure uses only stage-local state/history data; it does not
use endpoint-boundary source data, terminal-row replacement, projection, or a
target-direction oracle. \\
Algebraic health &
The solve keeps rank 132, residuals near the existing Newton tolerance, and
well-recorded conditioning diagnostics on the same smooth/sharp cases. \\
Trajectory evidence &
The candidate is promoted from local tangent probes to a nonlinear h-sweep over
$h=\{0.04,0.02,0.01\}$ against the $h=0.005$ reference, with CSV/JSON/plot
artifacts. \\
Order and closure &
The h-sweep must close raw terminal endpoint velocity while restoring the
paper-\tfe{} order target; the current source-free terminal-closed candidate
fails here because its smooth position order is 3.523. \\
Pipeline synchronization &
The summary JSON, full-\tfe{} gap ledger, repair spec, paper claim ledger, and
read-only validators must all move from the current
\nolinkurl{full_tfe_stage_replacement=false} boundary to the new accepted
claim. \\
\bottomrule
\end{tabular}
\end{table}

This target also explains why the recent negative audits are useful.  They do
not show that \tfe{} is impossible; they show which tempting shortcuts are not
the missing row.  Frozen row-space compression can balance values without
closing the tangent gap.  Near-terminal convex rows can approach the terminal
bridge but remain terminal-equivalent.  Current-pose and stage-difference
$C_v$ source-to-velocity lifts move source rows into generalized velocity
space, but still do not span the eight missing directions.  The next useful
implementation therefore has to change the weak-row formula itself, most
likely through a genuinely nonlinear recurrent history source law or a
stage-2 translational/angular coefficient-gradient term derived from the
lower-pair weak balance.

\section{Sparse Backend Status}

\plotfigure{\figpath/sparse_speed_gap.png}
{Sparse speed gap audit.  Row-VJP reduces the seed count but is not yet a
stable wall-clock improvement over dense \texttt{jacfwd}.}
{fig:sparse}

The row-colored VJP backend uses 60 row colors rather than the 90 column colors
of the dense-pattern coloring.  It preserves trajectory agreement, but the
current implementation is not a consistent runtime win over dense
\texttt{jacfwd}.  The sparse cost-model audit records the required row-runtime
reduction needed to match dense and to achieve a 10\% dense win.  This is an
engineering caveat, not a correctness caveat.

\section{Reproducibility Artifacts}

The primary generated artifacts are:
\begin{itemize}
  \item \artifact{v047\_cylindrical\_chain\_pipeline/run\_v047.py};
  \item \artifact{v047\_cylindrical\_chain\_pipeline/validate\_v047\_outputs.py};
  \item \artifact{v047\_cylindrical\_chain\_pipeline/validate\_full\_tfe\_gap.py};
  \item \artifact{v047\_cylindrical\_chain\_pipeline/validate\_full\_tfe\_repair\_spec.py};
  \item \artifact{v047\_cylindrical\_chain\_pipeline/FULL\_TFE\_REPLACEMENT\_GAP\_LEDGER.md};
  \item \artifact{v047\_cylindrical\_chain\_pipeline/FULL\_TFE\_REPAIR\_SPEC.md};
  \item \artifact{v047\_cylindrical\_chain\_pipeline/results/summary\_v047.json};
  \item \artifact{v047\_cylindrical\_chain\_pipeline/results/v047\_report.md};
  \item \artifact{ORDER\_PROOF\_LEDGER.md};
  \item \artifact{VALIDATION\_QUICKSTART.md};
  \item \artifact{paper\_v047\_cylindrical\_chain/CLAIM\_BOUNDARY.json};
  \item \artifact{paper\_v047\_cylindrical\_chain/REVIEWER\_CHECKLIST.md};
  \item the CSV and PNG files under \artifact{v047\_cylindrical\_chain\_pipeline/results/}.
\end{itemize}

The machine-readable claim boundary is \artifact{CLAIM\_BOUNDARY.json} in
this paper directory.  It records the accepted formal-order comparison, the
local paper-style comparator, the explicit non-claims, the open caveats, and
an \nolinkurl{evidence_files} manifest.  That manifest points to the summary
JSON, convergence CSV, ASME CSV rows, paper-\tfe{} formula-mapping CSV,
full-\tfe{} gap ledger, repair specification, paper PDF, claim ledger,
reviewer checklist, and Chinese status note.  The paper-claim validator parses
this JSON and checks those evidence paths before accepting the paper package.

At the time of this build, the validator expects the source-free closure, the
order/closure blend, and the velocity-compression probe artifacts through the
row-space compression, coefficient-derivative, state-feature coefficient-derivative,
and four-history source-law audits plus the JSON-only active-stage velocity
matrix-gradient localization.  The four-history source-law screen is negative
evidence: it keeps \texttt{terminal\_closed\_row\_count=0} while the best
smooth-order branch has min order \texttt{5.342}.  The rebuilt v047 artifact set passes the
v047 validator with 292 result files, 147 CSV files, and 120 PNG files; the
top-level pipeline validator also passes with 48 version directories, 695
result files, 286 CSV files, 202 PNG files, and 113 JSON files.

The full numerical regeneration is deliberately not the default paper build
path.  The latest recorded full v047 generation time is
$2386.76\,\mathrm{s}$, because \artifact{run\_v047.py} has accumulated the
complete historical audit suite.  For the four-example paper claim, the
authoritative evidence is the existing generated artifact set plus the
read-only validators.  A new full regeneration is needed only when numerical
CSV/JSON/PNG evidence changes.
Accordingly, \artifact{run\_v047.py} should be read as an audit harness and
artifact generator, not as a compact reference implementation of the accepted
integrator.  Its size reflects the accumulated negative full-\tfe{} screens and
historical diagnostics.  It is not necessary to execute that harness merely to
check the four-example claim or rebuild this paper.
The paper build and paper-claim check run from this paper directory; the v047
artifact checks run from the v047 pipeline directory.
The reviewer checklist in this directory gives the shortest read-only command
sequence for checking the accepted method claims, the diagnostic full-\tfe{}
boundary, and the explicit rule that \artifact{run\_v047.py} should not be run
for routine paper edits.
For full-\tfe{} blocker exploration, the preferred workflow is a targeted
\artifact{V047\_TARGET\_AUDIT} run that prints JSON for one local audit.  The
recent active-stage velocity matrix grid used this path, took 51.2 seconds,
and did not regenerate the full artifact set.  Later terminal-limit
Taylor-style probes used the same JSON-only pattern: an earlier pose+velocity
acceleration refinement covered 11 local rows in 30.7 seconds, found best
nonterminal row
\nolinkurl{stage02_convex_pose_velocity_extrapolate_0p01_0p02_posevelaccel_p0p0005_m0p0001_z}
at residual $1.0173\times10^{-5}$, and still found no terminal-bridge span.
The latest stage-2 source-to-velocity transport screen is the strongest
nonterminal transport clue so far: the scaled
\nolinkurl{stage02_convex_pose_velocity_0p00_sourceliftmatrixdiff_historydelta_m0p0001_z}
row reaches residual $1.251\times10^{-6}$, reduces the independent target rank
to six, and still has no span.  Adding a small angular/Lie pose shift does not
improve it: the best combined row reaches only $1.296\times10^{-6}$ and also
has no span.  A final three-row bare endpoint-pose velocity check finds that
\nolinkurl{stage02_convex_pose_velocity_0p00_z} spans the local terminal
bridge at roundoff, with residual $1.327\times10^{-15}$, but this is recorded
as a local row-space span rather than an accepted full-\tfe{} replacement: no
trajectory h-sweep, order proof, or artifact update is attached.  These probes
are intentionally diagnostic local evidence; they do not update the accepted
four-example artifact set.
A bounded JSON-only trajectory smoke now promotes that local span into the
nonlinear 132-row residual.  The \nolinkurl{stage02_convex_pose_velocity_0p00_z}
row closes terminal velocity to $8.124\times10^{-17}$ on a smooth three-$h$
short sweep, but its position/velocity orders are only $4.142/2.305$.  The
nearby nonterminal \nolinkurl{stage02_convex_pose_velocity_0p01_z} row
converges but leaves terminal velocity at $6.255\times10^{-6}$.  This is
trajectory-level diagnostic evidence, not an accepted full-\tfe{} replacement.
A non-projection extrapolation limit reaches terminal velocity
$1.254\times10^{-13}$ at
\nolinkurl{stage02_convex_pose_velocity_extrapolate_0p00001_0p00002_z}, but
the three-$h$ order remains $4.142/2.305$, so the terminal/order split is not
removed by taking the nonterminal points closer to the endpoint.

\section{Non-Claims and Limitations}

The accepted evidence is intentionally narrower than the complete research
program.  The paper therefore makes the following non-claims explicit.

\begin{table}[t]
\centering
\caption{Statements that are not claimed by the current v047 evidence.}
\label{tab:nonclaims}
\begin{tabular}{P{0.30\linewidth}P{0.58\linewidth}}
\toprule
Non-claim & Current reason \\
\midrule
The paper-\tfe{} method is fully accepted &
False.  The machine-readable gate remains
\nolinkurl{full_tfe_stage_replacement=false}; the accepted trajectory method
is the sixth-order \method{} path. \\
The four-example gate requires a new full regeneration &
False.  The four-example coverage rows are already accepted in the existing
artifact set and are checked by the lightweight four-example validator, but
the four-link and slider-crank rows are not dynamic-order proof rows. \\
\artifact{run\_v047.py} is a compact reference implementation &
False.  It is an accumulated audit harness and artifact generator containing
historical positive and negative screens. \\
Sparse AD is a speed win &
Not yet.  Row-VJP reduces seed count but is not a stable wall-clock
improvement over dense \texttt{jacfwd}. \\
Sharp Brown--McPhee friction is inexpensive at coarse steps &
Not yet.  Ultra refinement recovers high order, but the practical coarse
regime remains order-reduced and more expensive. \\
\bottomrule
\end{tabular}
\end{table}

\begin{table}[t]
\centering
\caption{Recommended reproducibility levels for this paper build.}
\label{tab:repro-levels}
\begin{tabular}{P{0.22\linewidth}P{0.42\linewidth}P{0.26\linewidth}}
\toprule
Level & Command or artifact & Purpose \\
\midrule
Paper build &
\artifact{cmd.exe /c latexmk -pdf -interaction=nonstopmode main.tex} &
Rebuilds the PDF from the current text and figure snapshot. \\
Paper package check &
\artifact{validate\_paper\_package.py} &
Runs the lightweight paper/four-example/artifact checks against the existing
PDF; pass \artifact{--latex} to attempt a rebuild from the wrapper. \\
Paper-claim check &
\artifact{validate\_paper\_claims.py} &
Checks that this paper's headline claims match the current summary, figures,
and PDF. \\
Four-example gate check &
\artifact{validate\_four\_asme\_minimal.py} &
Checks only the existing single, double, four-link, and slider-crank method
evidence. \\
Full-\tfe{} gap check &
\artifact{validate\_full\_tfe\_gap.py} &
Checks the open replacement boundary and verifies that terminal closure,
order/closure split, and compression blockers remain explicit. \\
Full-\tfe{} repair spec check &
\artifact{validate\_full\_tfe\_repair\_spec.py} &
Checks the next candidate's 132-row residual budget and the 16+8 lower-pair
implementation contract, including the bounded smoke entry. \\
Read-only artifact check &
\artifact{validate\_v047\_outputs.py} &
Checks the already-generated v047 CSV/JSON/PNG/report evidence without
rerunning the simulations. \\
Full numerical regeneration &
\artifact{run\_v047.py} &
Regenerates the full historical v047 audit suite; use this only after changing
the numerical method or artifact-producing code. \\
\bottomrule
\end{tabular}
\end{table}

\section{Remaining Gaps}

\begin{table}[t]
\centering
\caption{Current missing items after the accepted four-example method gate.}
\label{tab:remaining}
\begin{tabular}{p{0.27\linewidth}p{0.65\linewidth}}
\toprule
Gap & Current status \\
\midrule
Independent full-\tfe{} stage replacement &
Open.  Source-free terminal closure is available, but the accepted-order
nonlinear h-sweep is still missing.  The latest 36 target-free compression
rows, velocity-matrix, source/history, terminal-limit source-lift, curvature
source-lift scale, acceleration-bridge scale, generalized-acceleration shift,
pose-acceleration shift, tiny pose-acceleration scale, pose+velocity Taylor
shift, component-split pose-acceleration, component-split
velocity-acceleration, component-mixed pose/velocity Taylor, and Gauss
endpoint-pose probes all have zero spans.  The
closest terminal-limit source-lift family remains at
$5.298\times10^{-6}$ with source-curvature norm $1.076\times10^{-14}$, so the
next repair cannot be a simple source-lift scale tweak.  The strongest
nonterminal transport clue is the stage-2 source-to-velocity transport screen:
it reaches residual $1.251\times10^{-6}$ and lowers the independent target
rank to six, but it still has no span; the angular-pose combined follow-up
reaches only $1.296\times10^{-6}$ and remains no-span.  The bare
\nolinkurl{stage02_convex_pose_velocity_0p00_z} row spans locally at
$1.327\times10^{-15}$, confirming the endpoint limit of that family, but this
is not an accepted full-\tfe{} replacement.  The bounded nonlinear smoke closes
terminal velocity to $8.124\times10^{-17}$ for that row but gives only
$4.142/2.305$ position/velocity order; the nonterminal $0p01$ row leaves
terminal velocity at $6.255\times10^{-6}$.  Ultra-near linear extrapolation
and three-point quadratic nonterminal extrapolation both close terminal
velocity without projection, but both retain the same $4.142/2.305$ order
loss.  The diagnostic ledger also retains the missing-direction decomposition
and source-to-velocity evidence: the 34-row audit based on the
current-pose endpoint velocity matrix does not span; the matrix-difference
branch records 0.042593 before scale refinement, 0.009527 after scale refinement, and 0.009590
for the normalized-history source law.  The next repair must therefore derive a nonterminal endpoint-pose
predictor that closes terminal velocity on trajectory without order loss. \\
Sparse backend speed &
Open engineering caveat.  Row-VJP reduces seed count but is not yet a stable
runtime win over dense \texttt{jacfwd}. \\
Sharp-friction coarse regime &
Open practical caveat.  Deeper fixed refinement recovers high order, but the
coarse Brown--McPhee sharp case remains order-reduced and more expensive. \\
\bottomrule
\end{tabular}
\end{table}

\section{Conclusion}

The v047 evidence supports the intended conditional formal-order comparison.  The accepted
\method{} path follows the conditional sixth-order Gauss/FullVA theorem,
records smooth position/velocity orders 7.161/7.066 on the cylindrical
lower-pair benchmark, and has four-example coverage with dynamic-order
evidence on the single- and double-pendulum rows and closed-loop
residual/reaction evidence on the
four-link and slider-crank rows.  In the current artifact scope, this is the
result to compare against the local paper-style $m=3$
Lobatto-\tfe{} formula target with expected order five.  The sharp-friction
behavior is localized to a practical coarse-regime issue with high-order
recovery under deeper fixed refinement, not to a failure of the smooth
conditional sixth-order method claim.

The independent full-\tfe{} replacement remains open.  The source-free
final-stage velocity closure closes raw terminal endpoint velocity without
projection, endpoint source data, or terminal-row replacement, but its smooth
position order is too low for an accepted paper-\tfe{} replacement.  The newer
endpoint-pose velocity predictor smoke shows the same split at trajectory
level: the terminal-equivalent row closes velocity but loses order, while the
nearby nonterminal row does not close terminal velocity.  Ultra-near
non-projection extrapolation closes the terminal velocity tolerance at
$1.254\times10^{-13}$ but retains the same $4.142/2.305$ order loss; a
three-point quadratic nonterminal extrapolation similarly closes at
$2.029\times10^{-13}$ with finest-$h$ terminal velocity $8.254\times10^{-15}$
and the same order loss.  The
machine-readable
gate remains \nolinkurl{full_tfe_stage_replacement=false}.  The latest
compression
diagnostics rule out the tested target-free linear direction laws,
finite-difference second-differential non-final source features,
higher-order/nonlocal source features, one-step and two-step history-transport
features, recurrent curvature/bilinear history features,
source-active cross-Gram/Hadamard/shifted-commutator weak-row structures, and
target-free frozen row-space compressions.  The latest active-stage velocity
matrix-gradient localization improves the residual to 0.576548 but still has
rank-eight missing directions and no terminal-bridge span.  The best-row
decomposition puts 96.3\% of the remaining tangent energy in stage 2 and
identifies \nolinkurl{angular_velocity_w} as the dominant variable family.  A
simple angular-only mask follow-up reaches only 0.752997 and shifts the
remaining direction back to \nolinkurl{translation_velocity_v}; a direct
translation/angular cross-coupling follow-up reaches only 0.898812; an
endpoint-pose positive-Lagrange velocity predictor improves the local residual
to 0.117782 but still has rank eight/no span; a near-terminal convex predictor
has a terminal-equivalent local span but its best nonterminal row remains
rank eight/no span at 0.049317, 0.051890, and 0.052472 over
$h=0.04,0.02,0.01$; synchronized pose/velocity near-terminal rows reduce the
best local residual to 0.009512 but keep rank eight/no span; nonterminal
terminal-limit extrapolation reduces the residual to $4.699\times10^{-6}$ but
still has no span and is marked terminal-bridge-equivalent; direct
near-terminal slope rows reach only 0.677034; direct mean-acceleration bridge
correction worsens the residual to 0.076159 or above;
a direct near-terminal curvature screen reaches only 0.871228;
a source/history coefficient-feature matrix screen reaches only 0.898873;
a nonlinear recurrent curvature/history matrix-feature screen reaches only
0.898865;
a terminal-limit $C_v$ source-lift screen reaches only $5.298\times10^{-6}$;
a component-split velocity-acceleration Taylor screen reaches only
$2.138\times10^{-4}$ at ordinary scale and $9.8175\times10^{-6}$ at tiny
translational scale;
a component-mixed pose/velocity Taylor screen reaches only
$9.987\times10^{-6}$;
a stage-2-fixed/delta acceleration velocity-shift screen reaches only
$9.8175\times10^{-6}$;
a nonfinal terminal velocity/source predictor screen reaches only 0.923466;
a stage-2 source-to-velocity transport screen reaches $1.251\times10^{-6}$ but
no span, and an angular-pose combined follow-up reaches only
$1.296\times10^{-6}$ with no span; the bare
\nolinkurl{stage02_convex_pose_velocity_0p00_z} row spans locally at
$1.327\times10^{-15}$ and closes terminal velocity on a bounded trajectory
smoke, but the three-$h$ smooth orders are only $4.142/2.305$; the nonterminal
$0p01$ trajectory row leaves terminal velocity at $6.255\times10^{-6}$;
a Gauss endpoint-pose predictor screen reaches only 0.172659 and shifts the
remaining family to \nolinkurl{lie_position_u};
a bilinear
outer-product follow-up reaches only 0.898720; and a componentwise diagonal
follow-up reaches only 0.898530.  All likewise have no span.  The next
technical target is a revised analytical weak-row formula or richer nonlinear
recurrent history source law inserted into the nonlinear lower-pair h-sweep,
preserving terminal closure while restoring the expected paper-\tfe{} order.

\end{document}
